\documentclass{ieeeaccess}
% Official IEEE Access LaTeX template class, downloaded 2026-08-26
% from https://ieeeaccess.ieee.org/authors/submission-guidelines/.
\usepackage{cite}
\usepackage[utf8]{inputenc}
\usepackage{amsmath,amssymb}
\usepackage{booktabs}
\usepackage{tabularx}
\usepackage{array}
\usepackage{needspace}
\usepackage{float}
\usepackage{balance}
\newcolumntype{Y}{>{\raggedright\arraybackslash}X}
\usepackage{graphicx}
\usepackage{microtype}
\usepackage{algorithm}
\usepackage{algpseudocode}
\usepackage{xcolor}
\usepackage{xurl}
\usepackage[hidelinks]{hyperref}
\Urlmuskip=0mu plus 1mu
\emergencystretch=3em
\sloppy

\begin{document}
\history{}
\doi{}

\title{Confidence-Guided Selective LLM Re-Evaluation for Depression-Risk-Related Emotion Classification in Social Media Text}
% REQUIRED BEFORE SUBMISSION: replace the review placeholder with the final
% author list, numbered affiliations, corresponding author, and funding note.
\author{\uppercase{Author Details Omitted for Professor Review}}
\address[1]{Affiliation details will be inserted in the submission version.}
\tfootnote{Funding and acknowledgment information will be inserted in the submission version.}
\markboth
{Author \headeretal: Confidence-Guided Selective LLM Re-Evaluation}
{Author \headeretal: Confidence-Guided Selective LLM Re-Evaluation}
\corresp{Corresponding-author details will be inserted in the submission version.}

\begin{abstract}
Transformer classifiers can process social-media text at scale, but confidence-based selection does not ensure that a second model will correct the selected predictions. We evaluate a two-phase framework for three-class proxy emotion classification (Depression, Neutral, and Happy) that separates routing from LLM re-evaluation. Phase~1 uses temperature-calibrated DistilBERT predictions and a calibration-set risk--coverage criterion to route low-confidence inputs. In a three-seed comparison on the same 12,000 Reddit test posts, DistilBERT achieved 96.88 $\pm$ 0.04\% accuracy, versus 95.36 $\pm$ 0.08\% for Mistral~7B and 95.17 $\pm$ 0.07\% for Llama~2~7B frozen-backbone linear probes, with higher throughput and lower latency on a common A100 benchmark. The fixed operational checkpoint achieved 96.69\% accuracy and routed 171 posts (1.42\%). Re-evaluation using minimally sanitized original text yielded 96.67\% accuracy with Llama~2 Chain-of-Thought and 96.94\% with Llama~3 SELF-DISCOVER. On a 300-example synthetic Mixed Emotion stress test, accuracy increased from 81.33\% to 85.33\% and 87.33\%, respectively. Because prompt development reused the original routed Reddit cases, we additionally evaluated the frozen pipeline on 9,000 previously unused same-source posts. Routing 137 posts (1.52\%) raised accuracy from 96.73\% to 97.03\% with Llama~3, yielding 27 net corrections and a Holm-adjusted $p=0.000131$; the Llama~2 gain was not significant. These results support configuration-dependent correction with limited re-evaluation volume beyond the prompt-development sample. They concern proxy emotion labels, not clinical diagnoses; external validation and expert review remain necessary for broader use.
\end{abstract}

\begin{keywords}
Depression-risk-related emotion classification, social media text analysis, large language models, confidence-guided routing, Chain-of-Thought prompting, SELF-DISCOVER, mental health NLP, reasoning traceability
\end{keywords}

\titlepgskip=-21pt
\maketitle

\section{Introduction}

The World Health Organization estimates that approximately 332 million people live with depression \cite{ref1} and identifies it as a leading cause of ill health and disability \cite{ref2}. Depression and anxiety together account for an estimated US\$1 trillion in annual productivity losses \cite{ref3}; their prevalence also increased during the first year of the COVID-19 pandemic \cite{ref4}. This burden motivates research on emotional signals in everyday language, provided that computational text labels remain distinct from clinical assessment.

Validated questionnaires such as the Patient Health Questionnaire-9 (PHQ-9) support depression screening \cite{ref5}, while accessible, person-centred care remains a health-system priority \cite{ref6}. Social-media analysis serves a different purpose: it allows researchers to study linguistic and emotional patterns in non-clinical text \cite{ref7}. A post-level emotion prediction neither replaces those instruments nor establishes the author's clinical condition.

Domain-adapted models such as MentalBERT demonstrate the value of target-domain pretraining for mental health text classification \cite{ref8}. Nevertheless, emotion-fusion research identifies dataset quality and interpretability as persistent challenges \cite{ref9}. For user-level prediction, processing posts independently can also discard temporal and inter-post information \cite{ref10}. Our post-level task does not model a user's longitudinal history. Research on clinical explainability further emphasizes that useful explanations depend on the intended setting and the information clinicians need \cite{ref11}; a generated rationale alone does not meet those requirements.

Existing research establishes the value of transformer classification, confidence calibration, selective prediction, and LLM-generated explanations, but these components answer different questions and are not interchangeable. A calibrated classifier may identify difficult inputs without correcting them, while an LLM may generate a plausible rationale without improving the underlying decision. This leaves a practical evaluation gap: confidence-based case selection and reasoning-based correction should be connected in one fixed pipeline and assessed separately at the classifier, router, and re-evaluator levels.

To address this gap, we propose a confidence-guided two-phase framework for depression-risk-related proxy emotion classification that combines transformer-based classification with selective, rationale-supported LLM re-evaluation. The completed Phase~1 implementation uses DistilBERT to produce initial predictions and calibrated maximum-softmax-probability confidence scores. A dedicated calibration split is held out from model selection and used for temperature scaling, calibration assessment, and routing-threshold selection. Samples whose calibrated confidence scores fall below the selected threshold are routed to Phase~2. The threshold maximizes Phase~1 coverage subject to a one-sided upper bound on accepted-set selective risk. In Phase~2, low-confidence samples are re-analyzed through model-specific structured reasoning processes \cite{ref17}, \cite{ref18}. The resulting end-to-end evaluation reports not only final accuracy, but also whether routing concentrates errors and whether each re-evaluator corrects more routed errors than it introduces.

The evaluation addresses three research questions. \textbf{RQ1a} assesses the operational classifier's predictive and calibration performance; \textbf{RQ1b} compares the candidate classifiers under a shared data and search budget. \textbf{RQ2} asks whether the fixed routing policy concentrates Phase~1 errors while retaining high coverage. \textbf{RQ3} asks whether re-evaluation produces positive net corrections on Reddit posts and on a controlled synthetic stress test. Separating these questions allows a useful router and an ineffective re-evaluator to be identified within the same experiment.

\section{Study Contributions}

The study makes four concrete contributions:
\begin{itemize}
\item It operationalizes selective LLM re-evaluation as an end-to-end pipeline in which an efficient classifier handles high-confidence inputs and reasoning models receive only calibration-defined low-confidence inputs. The Phase~1 comparison supports operational model selection through predictive performance and measured inference throughput, latency, and memory on a common accelerator, rather than claiming exhaustive optimization of every 7B architecture.
\item It replaces a heuristic confidence cutoff with a reproducible calibration protocol: temperature scaling, a prespecified threshold grid, a one-sided selective-risk upper bound, and maximum-coverage selection among feasible thresholds. The selected policy is fixed before held-out and Phase~2 evaluation.
\item It evaluates routing quality separately from re-evaluation quality. Coverage, selective risk, error capture, corrected errors, introduced errors, and net corrections reveal whether the router finds difficult inputs and whether the LLM actually improves them. An additional 9,000-post same-source holdout tests the frozen pipeline beyond the sample used during prompt development.
\item It contributes a controlled Mixed Emotion stress test and case-level reasoning audit trail for studying posts whose surface cues and dominant emotional trajectory diverge. The stress test is explicitly separated from model training and threshold selection.
\end{itemize}

The contribution is the integration and evaluation protocol, not a new calibration estimator, language-model architecture, or reasoning method.

\section{Literature Review}

\subsection{Early Approaches: Machine Learning for Depression Detection}

Early approaches used linguistic and behavioral features with conventional classifiers. In a study of college-student essays, Rude et al. \cite{ref20} found more negative-emotion words and more frequent use of ``I'' among currently depressed than never-depressed participants. This finding concerns a specific population and writing task, rather than a universal diagnostic marker. De Choudhury et al. \cite{ref21} used linguistic style, engagement, and network features to predict depression-related outcomes from Twitter activity. Tsugawa et al. \cite{ref22} evaluated Twitter activity features with an SVM against questionnaire-based depression scores. CLPsych 2015 subsequently provided shared user-level classification tasks based on self-reported diagnoses and matched controls \cite{ref23}.

Feature-based approaches also included Shen et al.'s multimodal dictionary-learning model, which combined six groups of depression-related and behavioral features from Twitter \cite{ref29}. Tadesse et al. \cite{ref31} evaluated linguistic and topic-based features with several classifiers on 1,841 Reddit posts, comprising 1,293 depression-indicative and 548 comparison posts. These studies established useful predictive baselines without making their labels equivalent to clinical assessments. Wongkoblap et al. \cite{ref24} identified data-collection bias, consent, dataset quality, and methodological choices as continuing limitations of social-media mental-health research.

\subsection{Deep Learning for Contextualized Depression Detection}

Neural approaches learn text representations rather than relying exclusively on handcrafted features. Word embeddings provide dense representations that encode semantic relationships \cite{ref25}, \cite{ref26}. LSTMs were designed to learn dependencies over extended sequences \cite{ref27}, while CNNs operating on word vectors were evaluated for sentence classification \cite{ref28}. These are methodological foundations, not themselves evidence of depression-detection performance.

Orabi et al. (2018) \cite{ref30} compared several deep learning models for user-level depression detection on Twitter. Using a TF-IDF linear SVM as a baseline, they showed that CNN-based models, particularly those using their optimized embeddings, outperformed the baseline classifier. In their experiments, CNN models also achieved better performance than the BiLSTM model. The authors further noted that supervised deep neural approaches have seen limited adoption in this domain due to the difficulty of obtaining sufficiently annotated training data \cite{ref30}.

Shared evaluations also changed the prediction setting. The eRisk 2018 depression task released users' posts in chronological chunks and evaluated early decisions \cite{ref32}. This user-history setting differs from the independent post-level classification studied here; results across the two tasks are not directly comparable.

\subsection{Large Language Models for Mental Health Analysis}

Transformer encoders and generative LLMs serve different roles in mental-health text analysis. Encoders provide contextual representations for supervised classification \cite{ref8}, \cite{ref33}, \cite{ref34}, whereas generative models can produce labels together with textual responses. Their use in the same application domain does not make predictive accuracy and explanation quality interchangeable outcomes.

Zhang et al. \cite{ref33} studied transformer-based depression-risk detection using Sina Weibo posts. MentalBERT \cite{ref8} instead emphasizes domain-specific pretraining, while ensemble approaches combine transformer classifiers \cite{ref34}. These studies motivate contextual representations but do not determine whether a second-stage model should process every post or only a selected subset.

Shin et al. \cite{ref35} evaluated GPT-3.5 and GPT-4 on semistructured emotional diaries, using questionnaire-based assessments as reference measures. Their study supports investigating diary text for depression-related screening, but its participants, reference measures, and input format differ from subreddit-derived post labels.

Wang et al. \cite{ref43} used Reddit writings to predict questionnaire-related depression levels and generate explanations for individual questionnaire items. Lan et al.'s DORIS framework \cite{ref44} instead uses LLM-derived symptom annotations and temporal mood-course features with a gradient-boosting-tree classifier, alongside generated justifications. These are distinct ways of combining prediction and explanation, not evidence that a generated rationale necessarily corrects an initial label. Omar and Levkovich's review \cite{ref45} includes both encoder models, such as BERT and RoBERTa, and generative approaches; it identifies predictive potential while emphasizing further validation, privacy, and ethical safeguards. Its findings should not be read as established therapeutic efficacy of generative chatbots.

These approaches retain the construct-validity, sampling, and platform-dependence concerns documented in reviews \cite{ref36}, \cite{ref37}, \cite{ref46}. For the present task, a generated explanation is therefore an inspectable output, not evidence that the proxy label or the prediction is clinically valid.

\subsection{Confidence Calibration, Selective Prediction, and Cascaded Re-Evaluation}

Calibration addresses the correspondence between predicted confidence and observed correctness \cite{ref15}. Selective classification addresses a related but distinct problem: trading coverage for accepted-set risk by abstaining on selected inputs \cite{ref16}. Abstention does not itself calibrate probabilities or correct the rejected predictions. This distinction motivates separate evaluation of confidence quality, case selection, and downstream correction.

Here, rejected inputs are sent to an LLM rather than left unclassified. Temperature scaling \cite{ref15} and selective prediction \cite{ref16} motivate the routing design; CoT \cite{ref17} and SELF-DISCOVER \cite{ref18} motivate the adapted re-evaluation protocols described below. The empirical question is whether recombining their outputs improves the original classifier after newly introduced errors are counted.

\subsection{Dataset Considerations in Social Media-Based Mental Health Text Analysis}

Progress in AI-based depression-related text analysis has been strongly shaped by the availability and quality of datasets. Early studies often relied on relatively small, task-specific datasets collected from social media platforms where users publicly disclosed a diagnosis of depression. A representative example is the CLPsych 2015 Shared Task dataset, which comprised Twitter users who self-reported depression and demographically matched control users. This dataset served as an early benchmark for depression-related classification from social media text \cite{ref23}.

As research in this area expanded, social media corpora became increasingly important. For example, Tadesse et al. analyzed Reddit posts using natural language processing and machine learning techniques to identify depression-related content \cite{ref31}. Twitter was reported as the most frequently studied platform for social media-based depression-related analysis, while Reddit and Facebook were also used in this line of research \cite{ref36}.

Despite their usefulness, these datasets have important limitations. Many rely on self-disclosed diagnoses and indirectly constructed labels, raising concerns about construct validity and the operationalization of mental health status in social media-based research. Chancellor and De Choudhury identified inconsistencies in label operationalization, the lack of standardized evaluation practices, and the need for greater transparency in data curation and validation procedures in social media-based mental health research \cite{ref37}. More recent review evidence similarly emphasizes that social-media depression detection studies remain vulnerable to sampling bias, platform bias, language and demographic imbalance, inconsistent preprocessing, and incomplete reporting of model-development procedures \cite{ref46}. These concerns motivate the present study's explicit treatment of subreddit-derived labels as proxy emotion labels rather than clinical labels, as well as its reporting of data construction, filtering, calibration, and threshold-selection procedures.

Privacy and ethical concerns also remain central in mental health-related text analysis using social media data \cite{ref38}, \cite{ref39}. Because such data may contain sensitive personal information, responsible data handling practices and clearly defined research protocols are essential for ethical and rigorous research \cite{ref39}.

\section{Dataset and Preprocessing}

The primary Reddit corpus supplies model development and evaluation data. A separate synthetic dataset tests controlled affective scenarios and Neutral controls. Their construction and uses are described below.

\subsection{Dataset Construction}

The primary dataset was constructed from Reddit posts collected from subreddit communities used as class proxies. Depression and AnxietyDepression supplied the Depression group; technology, datascience, AskScienceDiscussion, and webdev supplied the Neutral group; and Positivity, happy, MadeMeSmile, UnexpectedlyWholesome, and CongratsLikeImFive supplied the Happy group. After cleaning, polarity filtering, and class balancing, the preprocessing artifact contained 133,660 rows for three-class emotion classification. Each row corresponds to a sampled Reddit post, and the post title and body were concatenated before cleaning. Because community membership supplies a proxy label rather than a clinical annotation, community-specific lexical cues and sampling choices are potential sources of construct bias.

The Depression class was intended to represent posts containing depression-related emotional expressions or distress-related language. The Neutral class represented posts with a relatively neutral affective tone, whereas the Happy class represented posts containing positive emotional expressions. These labels were used as proxy labels for depression-risk-related emotion classification and should not be interpreted as formal clinical diagnostic labels.

\subsection{Data Features}

For each Reddit post, both textual content and auxiliary metadata were retained during dataset construction. The Phase~1 modeling input was the cleaned title--body text (\path{Title_with_selftext_cleaned}), which was generated by concatenating the post title and body text. The retained original title and selftext were used only when a routed post entered Phase~2. Additional fields, including subreddit source, author-related metadata, adult-content flag, timestamp, polarity score, and class label, were used only to support preprocessing, sampling, linkage, and dataset organization. A complete description of the retained variables is provided in Appendix Table A1.

Class labels were numerically encoded for model training as Depression = 0, Neutral = 1, and Happy = 2. This encoding was used only for computational convenience and does not imply an ordinal relationship among the classes.

\subsection{Text Cleaning and Normalization}

A standardized preprocessing pipeline was applied to reduce noise and improve comparability across posts. The concatenated title and selftext fields were first converted to lowercase and tokenized. Non-alphabetic content, URLs, numbers, special characters, and common stop words were removed. Lemmatization was then applied to reduce inflected word forms to their base forms. For example, ``Sad'' was converted to ``sad,'' and ``running'' was lemmatized to ``run.''

In addition, a Reddit-specific stop-word list filtered platform-specific filler expressions. The cleaned field, \path{Title_with_selftext_cleaned}, was used for Phase~1 training, calibration, and prediction. Cleaning preserves the order of retained tokens but can remove negation, punctuation, and discourse markers. Phase~2 therefore used the minimally sanitized original text to retain these cues and sentence boundaries.

\subsection{Sentiment-Aware Data Sampling}

To support label consistency, a sentiment-aware filtering step was applied after text cleaning. The implementation passes the complete cleaned title--body string to TextBlob and reads its document-level polarity value in the range $[-1,1]$, as documented by TextBlob \cite{ref41}. Pattern provides related NLP and sentiment-analysis functionality \cite{ref42}; the expression below specifies the TextBlob call used here:

\begin{equation}
P_{post}=\mathrm{TextBlobPolarity}\!\left(\mathrm{Clean}(\mathrm{title}\oplus\mathrm{body})\right),
\label{eq:post-polarity}
\end{equation}

where $\oplus$ denotes title--body concatenation and $\mathrm{Clean}(\cdot)$ denotes the cleaning pipeline described above. Values near $-1$, $0$, and $+1$ indicate increasingly negative, neutral, and positive polarity, respectively. This expression matches the document-level software call used by the implementation; the study code does not compute a separate TextBlob score for every token and then average those scores.

Class-specific polarity thresholds were then applied to retain posts whose overall sentiment broadly aligned with the intended emotional label. Specifically, posts in the Depression class were retained if their polarity scores fell within [-1, 0.2], Neutral posts were retained within [-0.3, 0.3], and Happy posts were retained within [-0.2, 1]. Posts whose polarity scores were substantially inconsistent with their intended labels were excluded during dataset construction. This procedure was intended to reduce obvious sentiment-label mismatch while preserving the three-class emotion classification setting.

\subsection{Mixed Emotion Dataset}

In addition to the primary Reddit dataset, we constructed a supplementary synthetic Mixed Emotion Dataset to examine model behavior under emotionally ambiguous conditions. The final dataset contains 300 controlled examples, balanced across the three proxy emotion labels, with 100 Depression, 100 Neutral, and 100 Happy examples. This dataset was not used for Phase~1 model training, hyperparameter tuning, temperature scaling, or routing-threshold selection. Instead, it was designed as a targeted stress-test set for cases in which positive, neutral, and depression-related cues co-occur or shift across the text. In the evaluation workflow, the temperature and routing threshold selected on the Reddit calibration split are fixed and then applied directly to the Mixed Emotion Dataset; thresholds are not re-selected on the synthetic stress-test data.

The current dataset includes seven controlled scenario groups: blended-emotion co-occurrence (40), positive-to-distress shift (40), distress-to-recovery shift (40), neutral framing with subtle affect (40), conflicting cues with a dominant trajectory (40), clear technical or informational Neutral examples (75), and Neutral examples with mild factual or procedural ambiguity (25). Each example was assigned a target label according to the dominant overall emotional trajectory rather than isolated sentiment-bearing phrases. The two Neutral groups preserve the proxy-label definition while avoiding a stress set that is artificially dominated by either completely trivial or clinically suggestive Neutral content. This design is intended to test whether confidence-guided routing and reasoning-based re-evaluation can address cases that are plausibly more difficult for a single-stage classifier.

Synthetic examples were generated using controlled prompting and checked against inclusion and exclusion criteria. Kang et al. \cite{ref19} used LLM-generated clinical-interview summaries for depression-prediction training augmentation. Here, synthetic posts serve a different purpose: supplementary stress-test evaluation, not training augmentation or clinical validation. Inclusion required consistency with the assigned scenario and target label: mixed or shifting cues for the affective scenarios, and low-affect content for the Neutral controls. Exclusions covered explicit clinical diagnosis claims, treatment recommendations, crisis language, identifying information, and off-scenario content. Appendix Table~A5 records the generation protocol. These checks were not an independent expert-annotation study; results describe the final 300 synthetic examples rather than naturally occurring mixed-emotion performance.

\subsection{Data Partitioning and Model Inputs}

After preprocessing and label encoding, the primary Reddit dataset is divided into training, model-validation, threshold-calibration, and held-out test sets by independently shuffling and partitioning each class. The final protocol samples exactly 40,000 examples per class (120,000 in total) and assigns 28,000 examples per class to training and 4,000 examples per class to each of validation, calibration, and held-out testing. Thus, the final split sizes are 84,000 training, 12,000 validation, 12,000 calibration, and 12,000 test examples. This per-class construction avoids one-example rounding differences from successive proportion-based split operations.

The training partition supplies supervised parameter updates; model validation supplies hyperparameter, epoch, and checkpoint selection. The calibration partition supplies temperature fitting and threshold selection. Test labels are excluded from these Phase~1 operations. However, the routed Reddit test cases were reused during Phase~2 prompt development, as disclosed in the experimental limitations. Each classifier uses its own tokenizer on the cleaned title--body field.

Only the cleaned textual content was used as the direct Phase~1 classifier input. For routed Reddit examples, Phase~2 received the linked original title and selftext after minimal privacy-oriented sanitization. Auxiliary metadata, such as subreddit source, author-related metadata, timestamp, and adult-content flag, was not used as a direct predictive feature in either phase.

\subsection{Ethical Considerations}

Because this study analyzes depression-risk-related emotional signals from social media text, several ethical considerations are relevant to dataset construction and interpretation.

First, the dataset labels are proxy labels derived from social media context and sentiment-aware filtering, not diagnoses assigned by licensed mental health professionals. Therefore, the proposed framework should be interpreted as a research-oriented emotion classification system rather than a clinical diagnostic or screening tool.

Second, Reddit posts may contain sensitive personal expressions related to emotional distress or mental health. To reduce privacy risks, personally identifying information was not used as a modeling feature. Before routed original text was passed to Phase~2, URLs were replaced by \texttt{[URL]} and direct Reddit username patterns were replaced by \texttt{[USER]}; sentence order, punctuation, capitalization, and negation were otherwise preserved. Author-related metadata was used only for dataset organization and was not used for prediction. These measures are limited sanitization, not a guarantee of anonymity: retained verbatim passages may remain searchable. The analysis focused on aggregate model development and evaluation rather than individual-level identification or profiling.

Third, the Mixed Emotion Dataset was generated through controlled prompting with GPT-5 and retained as a versioned evaluation artifact. Because synthetic examples may reflect generation-model biases or stylistic patterns, this dataset was used only for stress-test evaluation and not for primary model training. Future work should validate the framework with naturally occurring and expert-reviewed datasets before considering any real-world decision-support use.

Finally, the framework was not designed to identify, monitor, or intervene with individual users. Any deployment beyond research use would require institutional review, domain-expert oversight, bias and fairness assessment, and compliance with relevant data governance requirements.

\section{Methodology}

The pipeline has two stages: an operational DistilBERT classifier \cite{ref12} assigns a label and calibrated confidence, and a separately configured LLM re-evaluates low-confidence inputs. Phase~2 uses either Llama~2-7B-Chat \cite{ref14} with Chain-of-Thought prompting \cite{ref17} or Llama~3-8B-Instruct \cite{ref40} with SELF-DISCOVER \cite{ref18}. The configurations are evaluated separately, not ensembled.

\subsection{Confidence-Guided Hybrid Model Architecture}

Development and inference are separate operations. Development fixes the classifier, calibration temperature, and routing threshold. At inference, predictions at or above the threshold retain their Phase~1 labels. Predictions below it receive a Phase~2 label and a stored rationale. Full-set evaluation then recombines the two paths, counting both corrected and newly introduced errors.

Figure~\ref{fig:architecture} summarizes the complete data-to-output workflow. It separates model development and calibration from held-out evaluation, shows the fixed accept-or-route decision, and makes explicit that the two Phase~2 configurations are evaluated separately rather than ensembled.

\begin{figure*}[t]
\centering
\includegraphics[width=\textwidth]{figures/architecture_figure_publication_final.pdf}
\caption{Confidence-guided two-phase pipeline. Phase~1 produces a calibrated confidence score using a temperature $T^*$ fitted on the calibration split and applies a fixed routing threshold $\tau^*$. Accepted predictions retain the Phase~1 label, while routed cases are re-evaluated by one separately configured Phase~2 reasoner and then recombined for end-to-end evaluation and audit.}
\label{fig:architecture}
\end{figure*}

\subsection{Phase 1: Initial Emotion Classification}

The Phase~1 comparison includes DistilBERT \cite{ref12}, Mistral~7B \cite{ref13}, and Llama~2~7B \cite{ref14} under common data partitions and a four-trial search budget. DistilBERT is fully fine-tuned, whereas the 7B backbones use frozen representations with trainable linear heads. Thus, this is a comparison of practical classifier configurations, not an isolated test of architecture or maximum attainable performance. Each selected configuration is evaluated with seeds 42, 43, and 44. Appendix Table~A1c gives the selected settings; only the fixed operational DistilBERT checkpoint supplies downstream routing scores.

\subsection{DistilBERT}

DistilBERT \cite{ref12} is a compact bidirectional encoder obtained by distilling BERT into a six-layer student. Its original pretraining combines language-modeling, teacher-distribution matching, and representation-alignment objectives, producing a model with about 40\% fewer parameters than BERT-base and a reported inference speedup of about 60\% while retaining much of the teacher's language-understanding performance. Those properties make it an appropriate candidate for the high-volume first stage of a selective pipeline, where every input must be classified before only a small subset is escalated. In this study, the released \texttt{distilbert-base-uncased} checkpoint is fully fine-tuned with supervised three-class cross-entropy; the historical distillation process is not repeated. Its logits are subsequently temperature-calibrated, and only this operational checkpoint supplies the confidence scores used for routing.

\subsection{Mistral}

Mistral~7B \cite{ref13} is a decoder-only transformer with grouped-query attention (GQA) and sliding-window attention (SWA). GQA reduces key--value memory requirements, while SWA restricts attention to a local window. Here, Mistral provides a larger pretrained representation for three-class classification. Its backbone remains frozen, and a bias-free linear head maps the 4,096-dimensional representation to three class logits. Only the head's 12,288 parameters are trained; no adapter is added.

\subsection{Llama 2}

Llama~2 \cite{ref14} is a family of decoder-only autoregressive transformers. Its 7B base checkpoint supplies a second larger representation family for Phase~1 and is distinct from the chat-tuned checkpoint used in Phase~2. Classification uses the same frozen-backbone, bias-free 12,288-parameter linear-head design as Mistral. The shared partitions and label mapping permit predictive comparison, while the differing adaptation scope precludes attributing differences to model size alone.

All three classifiers output class logits. Post-hoc calibration and threshold selection are performed only for the operational DistilBERT checkpoint; no calibrated routing result is claimed for the comparison models.

\subsection{Calibrated Confidence-Based Routing and Risk-Coverage Threshold Selection}

The fixed operational DistilBERT classifier supplies the logits used for calibration and routing. For an input text $x_i$, it outputs logits $\mathbf{z}_i$ over $K$ proxy emotion classes. The uncalibrated class probability for class $k$ is computed with the softmax function:

\begin{equation}
p_{i,k}=\frac{\exp(z_{i,k})}{\sum_{j=1}^{K}\exp(z_{i,j})}.
\end{equation}

The Phase 1 predicted label is defined as:

\begin{equation}
\hat{y}_i = \arg\max_k p_{i,k}.
\end{equation}

Because raw maximum softmax probability may be poorly calibrated \cite{ref15}, a dedicated calibration split $D_{cal}$ is used to fit a scalar temperature parameter $T$. Temperature-scaled probabilities are computed as:

\begin{equation}
q_{i,k}(T)=
\frac{\exp(z_{i,k}/T)}
{\sum_{j=1}^{K}\exp(z_{i,j}/T)}.
\end{equation}

The temperature is selected by minimizing negative log-likelihood on $D_{cal}$:

\begin{equation}
T^*=\arg\min_{T>0}
\frac{1}{|D_{cal}|}\sum_{i\in D_{cal}} -\log q_{i,y_i}(T).
\label{eq:temperature-nll}
\end{equation}

NLL evaluates the probability assigned to the true class and penalizes confident mistakes strongly. Two complementary diagnostics are also reported. The multiclass Brier score measures the squared distance between the complete probability vector and the one-hot target vector:

\begin{equation}
\mathrm{Brier}(T)=\frac{1}{|D_{cal}|}
\sum_{i\in D_{cal}}\sum_{k=1}^{K}
\left(q_{i,k}(T)-\mathbf{1}[y_i=k]\right)^2.
\label{eq:brier}
\end{equation}

Expected calibration error (ECE) partitions predictions into confidence bins $B_m$ and compares mean confidence with empirical accuracy within each bin:

\begin{equation}
\mathrm{ECE}(T)=\sum_{m=1}^{M}\frac{|B_m|}{|D_{cal}|}
\left|\mathrm{acc}(B_m)-\mathrm{conf}(B_m)\right|.
\label{eq:ece}
\end{equation}

NLL, Brier score, and ECE are complementary rather than interchangeable selection constraints. NLL is the optimization objective for $T^*$; Brier score checks the quality of all class probabilities; and ECE checks whether stated confidence agrees with observed correctness at the group level. Adaptive ECE repeats the last comparison with approximately equal-frequency bins to reduce sensitivity to sparse fixed-width bins. Lower values are better for all four diagnostics.

The primary routing confidence score is then the temperature-scaled maximum softmax probability:

\begin{equation}
c_i=\max_k q_{i,k}(T^*).
\end{equation}

For a candidate routing threshold $\tau$, the accepted set $A(\tau)$ and routed set $R(\tau)$ are defined as:

\begin{align}
A(\tau) &= \{i: c_i \geq \tau\}, \\
R(\tau) &= \{i: c_i < \tau\}.
\end{align}

Coverage and routing rate are computed as:

\begin{align}
\mathrm{Coverage}(\tau) &= \frac{|A(\tau)|}{N}, \\
\mathrm{RoutingRate}(\tau) &= \frac{|R(\tau)|}{N} \\
&= 1-\mathrm{Coverage}(\tau).
\end{align}

Selective risk is computed only on accepted Phase 1 predictions:

\begin{equation}
\mathrm{Risk}(\tau)=
\frac{1}{|A(\tau)|}
\sum_{i \in A(\tau)}
\mathbf{1}[\hat{y}_i \ne y_i].
\end{equation}

To make the routing rule more conservative on finite calibration samples, the implementation also computes a one-sided binomial upper confidence bound $\overline{R}_{\delta}(\tau)$ for the accepted-set error rate using confidence level $1-\delta$. The final threshold is selected on $D_{cal}$ by maximizing coverage while satisfying the target selective-risk constraint $\alpha$:

Let $e(\tau)=\sum_{i\in A(\tau)}\mathbf{1}[\hat{y}_i\ne y_i]$ denote the number of accepted-set errors and $n(\tau)=|A(\tau)|$ denote the number of accepted examples. For $e(\tau)<n(\tau)$, the one-sided Clopper--Pearson upper bound used by the implementation is:

\begin{equation}
\overline{R}_{\delta}(\tau)=
\mathrm{Beta}^{-1}\!\left(
1-\delta;\,
e(\tau)+1,\,
n(\tau)-e(\tau)
\right),
\label{eq:clopper-pearson-bound}
\end{equation}

where $\mathrm{Beta}^{-1}(\cdot;a,b)$ is the inverse cumulative distribution function of a Beta$(a,b)$ distribution. When every accepted prediction is incorrect, the bound is set to 1; empty accepted sets are ineligible. The tail probability $\delta$ specifies the nominal confidence level, whereas $\alpha$ specifies the acceptable risk budget. This is a calibration-set selection criterion. Because temperature fitting and threshold search reuse the calibration data, pointwise Clopper--Pearson bounds do not establish a simultaneous post-selection guarantee for the resulting policy. Nor does the criterion guarantee risk control under distribution shift.

\begin{align}
\tau^* &= \arg\max_{\tau \in \mathcal{T}}
\mathrm{Coverage}(\tau) \\
\text{s.t.}\quad
\overline{R}_{\delta}(\tau) &\leq \alpha, \\
|A(\tau)| &\geq m_{\min}.
\end{align}

Here, $m_{\min}$ is a minimum accepted-sample requirement used to avoid selecting thresholds supported by too few calibration examples. The implementation sets $m_{\min}=\lceil0.10|D_{cal}|\rceil=1{,}200$. The selected $\tau=0.70$ accepted 11,830 calibration examples, so this safeguard was satisfied but did not determine the selected operating point. In the final DistilBERT operating policy, candidates are swept from 0.70 to 1.00 in increments of 0.01. This lower bound was chosen before held-out testing to avoid treating a weak three-class majority probability near 0.50 as sufficiently reliable for direct acceptance. Unless otherwise stated, the primary operating point uses temperature-scaled maximum softmax probability, a target selective risk of $\alpha=0.05$, a one-sided risk confidence level of $1-\delta=0.95$, the acceptance rule $c_i \geq \tau$, and the routing rule $c_i < \tau$. The risk constraint is a feasibility condition: it excludes candidate thresholds whose accepted-set upper risk exceeds $\alpha$. Among the feasible candidates in the prespecified 0.70--1.00 grid, the candidate with the greatest Phase~1 coverage is selected. Thus, the completed DistilBERT value $\tau^*=0.70$ reflects both the prespecified conservative operating range and calibration-set risk--coverage evaluation; it was not optimized using held-out test results or Phase~2 outcomes. The 7B models serve as Phase~1 predictive comparators only; they do not define alternative routing policies in the reported end-to-end experiment.

The threshold-selection output records accepted accuracy, routed Phase~1 errors, error capture, routing precision, and proxy-Depression false-negative risk among accepted predictions. Additional confidence diagnostics include expected calibration error (ECE), adaptive ECE, Brier score, and negative log-likelihood before and after temperature scaling. Auxiliary analyses compare raw MSP, temperature-scaled MSP, entropy-based certainty, and probability margin; estimate bootstrap confidence intervals for selective metrics; assess threshold stability through calibration-set resampling; and extract high-confidence accepted errors for qualitative auditing.

\begin{algorithm}[t]
\caption{Calibrated confidence-guided routing}
\label{alg:confidence-routing}
\footnotesize
\begin{algorithmic}[1]
\State \textbf{Input:} $D_{val}$, $D_{cal}$, classifier $f_{\theta}$
\State \textbf{Input:} $\alpha$, $1-\delta$, grid $\mathcal{T}$, minimum count $m_{\min}$
\State \textbf{Output:} temperature $T^*$ and threshold $\tau^*$
\State Select best checkpoint using $D_{val}$
\State Fit $T^*$ on $D_{cal}$ by minimizing NLL
\State Compute calibrated confidences $c_i$ on $D_{cal}$
\State Sweep the prespecified calibrated-confidence candidate grid
\For{each $\tau\in\mathcal{T}$}
    \State Form $A(\tau)$ and $R(\tau)$ from $c_i$
    \State Compute coverage, routing rate, and risk bound
\EndFor
\State Retain candidates with $\overline{R}_{\delta}(\tau)\leq\alpha$ and $|A(\tau)|\geq m_{\min}$
\State If none remain, report infeasibility and stop
\State Choose a retained $\tau$ with maximum coverage
\State Fix $T^*$ and $\tau^*$ before held-out testing
\end{algorithmic}
\end{algorithm}

\begin{algorithm}[t]
\caption{Two-phase inference with fixed calibrated routing policy}
\label{alg:two-phase-inference}
\footnotesize
\begin{algorithmic}[1]
\State \textbf{Input:} post record $x$, classifier $f_{\theta}$, fixed $T^*$ and $\tau^*$
\State \textbf{Input:} Phase~2 reasoner $g$ (Llama~2 CoT or Llama~3 SELF-DISCOVER)
\State Obtain the Phase~1 input $x^{(1)}$ from record $x$
\State Compute $\mathbf{z}=f_{\theta}(x^{(1)})$ and probabilities $q_k(T^*)$
\State Set $\hat{y}^{(1)}=\arg\max_k q_k(T^*)$ and $c=\max_k q_k(T^*)$
\If{$c \geq \tau^*$}
    \State \textbf{return} $\hat{y}^{(1)}$ as the accepted Phase~1 final label
\Else
    \State Retrieve the minimally sanitized original text $x^{(2)}$
    \State Generate rationale and terminal label with $g(x^{(2)},\hat{y}^{(1)})$
    \State Parse exactly one label from \{Depression, Neutral, Happy\}
    \State \textbf{return} the parsed Phase~2 final label
\EndIf
\end{algorithmic}
\end{algorithm}

\begin{algorithm}[t]
\caption{Audited end-to-end evaluation of routed predictions}
\label{alg:end-to-end-audit}
\footnotesize
\begin{algorithmic}[1]
\State \textbf{Input:} reference labels $y_i$, Phase~1 labels $\hat{y}^{(1)}_i$, route indicators $r_i$
\State \textbf{Input:} parsed Phase~2 labels $\hat{y}^{(2)}_i$ for routed inputs
\For{each example $i$}
    \State $\hat{y}^{(E)}_i \gets \hat{y}^{(2)}_i$ if $r_i=1$; otherwise $\hat{y}^{(E)}_i \gets \hat{y}^{(1)}_i$
\EndFor
\State $C\gets\sum_i\mathbf{1}[r_i=1,\hat{y}^{(1)}_i\ne y_i,\hat{y}^{(E)}_i=y_i]$
\State $I\gets\sum_i\mathbf{1}[r_i=1,\hat{y}^{(1)}_i=y_i,\hat{y}^{(E)}_i\ne y_i]$
\State Report corrected errors $C$, introduced errors $I$, net corrections $C-I$
\State Report end-to-end metrics over all examples using $\hat{y}^{(E)}$
\end{algorithmic}
\end{algorithm}

Algorithm~\ref{alg:end-to-end-audit} formalizes the evaluation boundary between routed-subset behavior and full-set performance. In particular, a Phase~2 model is not credited merely for correcting selected errors: any newly introduced errors are deducted, and the final accuracy is recomputed over the complete evaluation set after accepted Phase~1 labels and routed Phase~2 labels are recombined.

The implementation records an explicit infeasibility status if no candidate satisfies the risk constraint. A minimum-risk fallback exists only for notebook smoke tests and was not invoked in the reported paper-scale experiment; no fallback threshold is treated as satisfying the risk-control constraint. The selected threshold, confidence method, temperature, risk-control method, candidate-generation rule, acceptance boundary, tie-breaking rule, split ratios, random seed, feasibility status, and final selective metrics are saved in \texttt{threshold\_provenance.json}. This provenance records the fixed operating decision for reproducibility; it does not by itself establish independence from earlier development choices.

\subsection{Relationship to Mixed-Emotion Evaluation}

The routing mechanism is entirely confidence-based during inference. It does not directly detect blended emotion or sentiment shift. Instead, the Mixed Emotion Dataset serves a complementary evaluation role by stress-testing whether the Phase 1 classifier and the Phase 2 reasoning stage behave more reliably on examples where the one-label-per-input assumption is more difficult. This separation keeps the routing rule reproducible while allowing the evaluation to probe emotionally complex cases.

\subsection{Phase 2: Reasoning-Based Re-Evaluation with Llama 2 and Llama 3}

Phase~2 adds a reasoning step to re-evaluate and explain predictions that were flagged as uncertain or potentially incorrect in Phase~1. LLMs are deployed in a zero-shot prompting setup to perform this secondary analysis. Specifically, Llama~2-7B-Chat, referred to as Llama~2, is paired with Chain-of-Thought prompting \cite{ref17}, and Llama~3-8B-Instruct, referred to as Llama~3, is paired with SELF-DISCOVER reasoning \cite{ref18}. These models are not fine-tuned Phase~1 classifiers. Instead, they serve as secondary evaluators that consider the content of a post along with the initial model outcome, then provide a rationale and, when appropriate, a revised prediction. To preserve the evidence needed for this task, routed Reddit posts are linked back to their original title and selftext. The reasoning input replaces URLs and direct username patterns but retains sentence order, punctuation, capitalization, negation, and temporal transitions. The two prompting protocols use different intermediate reasoning formats, but both are constrained to end with exactly one explicit final label in the same three-class space: \texttt{Final label: Depression}, \texttt{Final label: Neutral}, or \texttt{Final label: Happy}. Thus, Phase~2 comparisons are performed at the final-label decision level rather than at the raw-output level. End-to-end correction results are reported only after the routed-sample runs are complete.

\subsection{Llama 2 with Chain-of-Thought Prompting}

Llama 2-7B-Chat, a chat-optimized variant of Llama 2, was used to produce a structured re-evaluation record for routed examples. Our zero-shot, multistage protocol follows the CoT idea of eliciting intermediate reasoning, rather than reproducing the few-shot demonstrations evaluated by Wei et al. \cite{ref17}. It elicits an independent assessment, comparison with Phase~1, and a terminal class label. The retained text is inspectable, but its presence alone is not evidence that the rationale is faithful or clinically valid.

For the final Llama~2 configuration, the CoT protocol first requires an independent, text-grounded assessment before revealing the Phase~1 label. It then compares the two decisions and ends with one exact label from Depression, Neutral, and Happy. The protocol explicitly avoids numerical emotion percentages, because a largest percentage component can conflict with the qualitative final decision. The complete reported prompting structure is detailed in Appendix Table A2.

The revised CoT protocol preserves a text-grounded rationale while avoiding the conflicting numerical-breakdown rule used in earlier exploratory prompting. Phase~2 output quality and correction behavior were evaluated empirically on the fixed routed subsets rather than inferred from prompt design alone.

\subsection{Llama 3 with SELF-DISCOVER}

The protocol is summarized visually in Appendix Figure~C1; the figure is provided as a reference aid rather than as a new model contribution.

We use an input-specific adaptation of SELF-DISCOVER \cite{ref18} with Llama~3-8B-Instruct, abbreviated as Llama~3 SELF-DISCOVER below. The original method discovers a task-level plan from unlabeled task examples and reuses it across instances. Our implementation instead constructs and stores a plan for each routed post. Its prompt contains the 39-module reasoning pool: prompt-level instructions, not neural-network layers or separately trained components. For each post, \texttt{SELECT} identifies potentially useful instructions, \texttt{ADAPT} tailors them to the classification problem, and \texttt{IMPLEMENT} composes a structured plan, generally represented in JSON. The model then applies the plan and returns an assessment ending in one permitted terminal label. Appendix Tables~A3--A3c document this adaptation and its prompt-development diagnostics.

For the observed case \texttt{MEV2\_DEP\_022}, Phase~1 predicted Happy for a post that begins with encouraging feedback but ends with unresolved sadness. The stored \texttt{SELECT} trace names Critical Thinking, Textual Analysis, Risk Analysis, and Reflective Thinking. The subsequent artifacts adapt these instructions and organize them into a JSON plan. The terminal response cites the later sadness and returns \texttt{Final label: Depression}. Appendix Table~A4b provides excerpts. This record shows the model's generated justification; it does not establish which cues caused either model's prediction.

The module references in these fields are generated reasoning traces. They may paraphrase or group the canonical module wording, so module identifiers are retained for auditability but are not themselves treated as quantitative predictions. Quantitative evaluation uses only the parsed terminal label, whereas corrected, introduced, and net errors determine whether the longer reasoning process improves classification.

\subsection{Model-Specific Prompt Protocols and Final Input Policy}

Exploratory fixed-subset diagnostics showed that a shared cross-model prompt did not benefit both reasoners. The final Llama~2 protocol therefore uses a frozen CoT prompt, which requires independent assessment before Phase~1-label comparison and one exact terminal label, while Llama~3 uses the SELF-DISCOVER protocol. These model-specific protocols were frozen before the final original-text evaluation. Appendix Table~A3c reports the earlier cleaned-input development diagnostics separately from the final end-to-end results so that prompt development is not confused with the final input-policy evaluation.

For the final Reddit run, the 171 routed records were linked to retained original title--selftext pairs by normalized exact matching of the cleaned Phase~1 field and proxy label. All 171 records were matched, no key mapped to conflicting original texts, and only one post exceeded the 6,000-character reasoning limit; for that post, the beginning and ending were preserved with an explicit middle-omission marker. This linkage audit was completed before LLM loading and is exported with the result package. The final policy therefore changes only the text representation supplied to Phase~2, not the Phase~1 model, temperature, threshold, routed IDs, LLM checkpoints, prompts, or generation settings.

\subsection{Operational Inference Paths}

Only routed inputs undergo Phase~2 inference. Accepted labels remain unchanged even if they are wrong; routed labels are replaced by the parsed terminal output even if that replacement introduces an error. The final runs produced valid labels for all routed records. Reference labels are used for evaluation, not provided to the reasoner. Selecting fewer inputs reduces the number of posts re-evaluated, but does not by itself quantify latency or monetary savings.

\subsection{Reasoning Outputs as Audit Artifacts}

A key benefit of the hybrid approach is that it adds an inspectable rationale layer to routed predictions. For each routed input, the final output is accompanied by a natural-language explanation generated by Llama~2 or Llama~3 that identifies the textual cues used to support the terminal label. This improves traceability, although a generated rationale should not be assumed to be faithful or clinically valid without independent review.

For example, if a routed post is assigned Depression, the explanation can identify expressions of persistent sadness, hopelessness, or unresolved emotional decline and contrast them with any isolated positive cues. Domain experts can then examine whether the cited evidence supports the final class. The explanation is therefore treated as an audit artifact, not as proof that the prediction is correct. Retaining the full reasoning fields, parsed label, initial prediction, and reference label enables both quantitative error accounting and qualitative review of failure modes.

\section{Experimental Setup}

\subsection{Computing Environment}

The experiments were implemented in Python with PyTorch, Hugging Face Transformers, Datasets, scikit-learn, SciPy, and pandas; W\&B recorded the model-selection sweeps. Phase~1 used \texttt{distilbert-base-uncased} as the operational checkpoint and prespecified 7B base checkpoints for the Mistral and Llama~2 comparison classifiers. Routed Phase~2 inference used \texttt{NousResearch/Llama-2-7b-chat-hf} and \texttt{NousResearch/Meta-Llama-3-8B-Instruct}, loaded in 4-bit form with bitsandbytes in GPU-backed Google Colab sessions. Llama~2 CoT generation used a 256-token limit with sampling temperature 0.6 and top-$p$ 0.9; the Llama~3 SELF-DISCOVER structure stages used a 1,024-token limit and the terminal answer was decoded deterministically. The repository notebooks print installed package versions at runtime and save prediction-level outputs after every routed row.

The three Phase~1 comparison configurations were trained and profiled on NVIDIA A100-SXM4-40GB accelerators in the same software environment. The inference benchmark used fp16, a 256-token cap, and batch size 32 over the 12,000-post held-out split; single-post latency was measured separately on 512 posts. Three warm-up batches were discarded, and each benchmark was repeated three times. CUDA synchronization bracketed timed regions, which included host-to-device transfers. NVML sampled GPU utilization and board power at 1~Hz. Appendix Table~A1d reports the measurement protocol and stage-level resource use. The additional 9,000-post evaluation ran on an NVIDIA L4 using the model-specific inference protocols described above.

\subsection{Model Selection and Calibration Protocol}

Model selection uses the training and validation partitions described above. Temperature fitting and threshold selection use the separate calibration partition after the operational checkpoint is fixed. Neither operation updates the classifier weights. Calibration diagnostics are measured on the temperature-fitting partition and therefore describe calibration fit, not independent evidence of calibration generalization. Routing performance is evaluated separately on the test partition.

\subsection{Confidence Analysis and Provenance Outputs}

For reproducibility, the advanced confidence workflow exports temperature-scaling metadata, calibration-quality metrics, threshold sweeps for each confidence score, bootstrap confidence intervals, threshold-stability summaries, high-confidence accepted errors, per-class selective metrics, fixed-threshold Mixed Emotion Dataset results, and Phase~2 end-to-end summaries. The results in this manuscript were derived from the saved \texttt{threshold\_provenance.json} file and prediction-level CSV files rather than manually copied notebook output. Phase~2 evaluation reports both corrected Phase~1 errors and introduced errors because a reasoner can improve some routed cases while worsening others.

\subsection{Hyperparameter Tuning and Training Optimization}

Each W\&B Bayesian sweep comprised four trials maximizing validation macro F1. The shared candidates were batch sizes 16, 32, and 64; training budgets of two or three epochs; and weight decays $10^{-2}$, $10^{-3}$, and $10^{-4}$. Learning-rate ranges differed for encoder fine-tuning and frozen-backbone probing. The selected configuration was rerun with seeds 42, 43, and 44 using the training partition, with validation macro F1 determining the retained epoch. A selected three-epoch budget therefore does not demonstrate convergence beyond that search limit. Appendix Table~A1c reports the sweep-selected settings and, in its note, the fixed operational checkpoint used for downstream analyses.

\subsection{Loss Function and Optimizer}

The classifiers minimize three-class cross-entropy using AdamW and gradient clipping. Updates apply to all DistilBERT parameters or to the linear head alone for the frozen 7B models.

\subsection{Evaluation Metrics}

The metrics were organized by the research question they answer rather than treated as interchangeable indicators. For RQ1a, predictive performance was evaluated with accuracy and class-wise precision, recall, and F1, with unweighted macro averages reported across Depression, Neutral, and Happy; calibration was evaluated with NLL, Brier score, ECE, and adaptive ECE. For RQ1b, matched classifier performance is reported as mean $\pm$ standard deviation across three seeds and interpreted together with model parameter scale, trainable scope, and the common-hardware inference benchmark. Timing variability is reported separately as mean $\pm$ SD across three measurement repetitions. For RQ2, routing was evaluated with coverage, routing rate, accepted-set risk and its one-sided upper bound, error capture, and routing precision. For RQ3, the primary end-to-end outcome was the example-level change in correctness after accepted Phase~1 predictions were recombined with routed Phase~2 labels. This change was summarized by full-set accuracy change and by corrected, introduced, and net errors. Calibration and routing measures are supporting diagnostics: they establish whether the confidence policy is credible and whether it concentrates difficult cases, but they do not by themselves establish that Phase~2 improves the final system.

Phase~1 and end-to-end correctness are paired for each example. Two-sided exact McNemar tests use the discordant counts: corrected Phase~1 errors and newly introduced errors. For the original Reddit and Mixed Emotion evaluations, paired percentile-bootstrap 95\% confidence intervals use 50,000 row-level resamples and seed 20260813, with Holm adjustment across the four dataset--reasoner comparisons. The additional holdout uses 10,000 paired correctness resamples and seed 20260911, with its two prespecified comparisons against Phase~1 treated as a separate Holm family. These analyses characterize the observed prediction pairs conditional on the fitted models and selected prompts; they do not account for prompt selection, repeated LLM generation, or dependence between posts from the same author. The three-seed intervals in the classifier comparison instead describe variation across training runs on a common test set.

Residual accepted-set errors were audited separately using the fixed 12,000-example held-out protocol and the calibration-selected threshold. The audit reports accepted selective risk, reference-to-prediction transitions, confidence-band error rates, and Depression false-negative risk among accepted reference-Depression posts. For all 310 accepted errors, the stored model input was linked back to the retained Reddit title and selftext by normalized exact matching against \texttt{title\_with\_selftext\_cleaned}; all 310 rows were matched. Six high-confidence cases were then selected because the original post itself provides comparatively clear evidence for the stored proxy label and exposes a distinct error pattern. These cases illustrate failure modes rather than estimate their prevalence. Appendix Tables~A4e--A4f report the reference label, prediction, calibrated confidence, minimally sanitized original-post excerpt, and text-grounded audit rationale. Qualitative review does not overwrite any reference label or alter any reported metric.

\subsection{Additional Same-Source Holdout Protocol}
\label{sec:additional-holdout-protocol}

After freezing the model-specific prompts and input policy, we formed an additional balanced holdout of 9,000 posts from unused records in the same source corpus. We excluded text matching any post in the previously used 120,000-post sample after Unicode NFKC normalization, case folding, and whitespace normalization, and removed within-pool normalized-text duplicates. Sampling used seed 20260911 and retained 3,000 posts per class. No selected text overlapped the original sample under this matching rule. This is a same-source holdout, not an author-disjoint, temporal, or external-domain test.

The saved operational DistilBERT checkpoint, temperature $T^*=1.7706$, threshold $\tau^*=0.70$, and model-specific prompt and input policies were retained without fitting or selection on these 9,000 posts. Both reasoners processed the same routed IDs in separate configurations. Accepted predictions remained unchanged, and any unparseable routed output would be counted as incorrect rather than silently replaced by Phase~1. Evaluation was a single fixed run per configuration, not a new three-seed training comparison. Detailed execution settings are retained in the reproducibility records.

\section{Experimental Results}

\subsection{Phase 1 Performance on Reddit Data}

The matched Phase~1 comparison used the same class-balanced 120,000-post sample, prespecified 70/10/10/10 partitions, 12,000-example held-out test set, maximum input length of 256 tokens, four-trial Bayesian search budget, and random seeds 42, 43, and 44 for all three architectures. Table~\ref{tab:phase1-performance} reports predictive performance and inference efficiency for these configurations; training configurations are documented in Appendix Table~A1c.

\begin{table*}[t]
\centering
\caption{Matched Phase~1 predictive performance and inference efficiency}
\label{tab:phase1-performance}
\footnotesize
\setlength{\tabcolsep}{4.2pt}
\renewcommand{\arraystretch}{1.16}
\textbf{(a) Predictive performance on the common 12,000-post test set}\par\smallskip
\begin{tabularx}{\textwidth}{@{}l *{4}{>{\centering\arraybackslash}X}@{}}
\toprule
Model & Accuracy, mean $\pm$ SD & Accuracy 95\% CI & Macro F1, mean $\pm$ SD & Macro F1 95\% CI \\
\midrule
DistilBERT & 96.88 $\pm$ 0.04\% & [96.79, 96.98]\% & 96.88 $\pm$ 0.04\% & [96.79, 96.98]\% \\
Mistral~7B & 95.36 $\pm$ 0.08\% & [95.15, 95.57]\% & 95.36 $\pm$ 0.08\% & [95.16, 95.56]\% \\
Llama~2~7B & 95.17 $\pm$ 0.07\% & [95.00, 95.34]\% & 95.17 $\pm$ 0.07\% & [94.99, 95.34]\% \\
\bottomrule
\end{tabularx}
\vspace{2pt}
\parbox{\textwidth}{\footnotesize \textit{Note.} Panel (a) reports percentages over training seeds 42, 43, and 44. Confidence intervals are two-sided Student-$t$ intervals across the three seed-level results. The comparison evaluates a bounded operational screening protocol; it does not estimate fully tuned 7B performance.}
\par\medskip
\textbf{(b) Inference on NVIDIA A100-SXM4-40GB}\par\smallskip
\begin{tabularx}{\textwidth}{@{}l >{\centering\arraybackslash}X >{\centering\arraybackslash}X >{\centering\arraybackslash}X@{}}
\toprule
Model & Throughput (posts/s) & Single-post p50 (ms) & Peak memory (GB) \\
\midrule
DistilBERT & 4,959.5 $\pm$ 15.5 & 4.80 $\pm$ 0.03 & 0.29 \\
Mistral~7B & 51.7 $\pm$ 0.0 & 36.62 $\pm$ 0.06 & 15.22 \\
Llama~2~7B & 55.5 $\pm$ 0.0 & 34.05 $\pm$ 0.10 & 14.05 \\
\bottomrule
\end{tabularx}
\par\smallskip
\parbox{\textwidth}{\footnotesize \textit{Note.} Panel (b) reports mean $\pm$ SD across three timing repetitions, distinct from the training seeds in panel (a). Throughput uses batch size 32 and all 12,000 test posts; p50 is the median latency for 512 posts processed individually. All models use fp16 and a 256-token cap. Memory refers to batched inference. An SD displayed as 0.0 reflects rounding. Appendix Table~A1d provides p95 latency and the full protocol.}
\end{table*}

Figure~\ref{fig:phase1-model-comparison} visualizes the same observed three-seed results and their 95\% confidence intervals. DistilBERT achieved 96.88 $\pm$ 0.04\% accuracy, compared with 95.36 $\pm$ 0.08\% for Mistral and 95.17 $\pm$ 0.07\% for Llama~2. Macro F1 showed the same pattern. Under this fixed comparison protocol, DistilBERT exceeded Mistral and Llama~2 by 1.52 and 1.71 percentage points in mean accuracy, respectively. The two probes differed by 0.19 points; their relative ordering is not the basis of the operational choice. 

\begin{figure*}[t]
\centering
\includegraphics[width=0.88\textwidth]{figures/figure_e_phase1_model_comparison.pdf}
\caption{Observed Phase~1 classifier performance on the common 12,000-example held-out Reddit test set. Columns show three-seed means and whiskers show Student-$t$ 95\% confidence intervals for accuracy and macro F1. These intervals describe training-run variation, not repeated timing measurements. The vertical axis is truncated to make the between-model differences legible.}
\label{fig:phase1-model-comparison}
\end{figure*}

DistilBERT had the highest mean score and approximately 67 million parameters, compared with about seven billion for each comparator. These results support retaining it for the tested pipeline, without establishing superiority over fully adapted 7B models. Depression--Happy confusions accounted for 72.8\% of DistilBERT errors, 75.6\% of Llama~2 errors, and 74.3\% of Mistral errors across pooled seed predictions; Neutral F1 exceeded 98\% for all three. Pooling summarizes error composition, not 36,000 independent test examples. Small probe seed deviations also reflect the frozen backbone and should not be interpreted as equivalent training variability across adaptation methods.

The matched comparison and the downstream experiment serve different purposes. The former summarizes three-seed classifier performance; the latter retains the independently fixed operational DistilBERT checkpoint with 96.69\% accuracy. Its predictions, temperature, and routed IDs were not replaced after the comparison. All calibration, routing, and end-to-end results refer to that single operational checkpoint.

\textit{Phase~1 inference efficiency.} Table~\ref{tab:phase1-performance}(b) adds a measured computational basis for this choice. On the common A100 benchmark, DistilBERT processed 4,959.5 posts/s, giving 89.4- and 95.9-fold higher batched throughput than the Llama~2 and Mistral probes. Its median single-post latency was 4.80~ms, compared with 34.05 and 36.62~ms, and its peak inference memory was 0.29~GB versus 14.05 and 15.22~GB. These are classifier inference measurements, not speedups of the complete two-phase system. Because Phase~1 processes every input, its low per-post cost complements selective use of the more expensive reasoners. The benchmark profiles the matched comparison configurations; it does not replace or remeasure the fixed downstream checkpoint.

The one-off classifier build costs were much closer: 0.59, 0.68, and 0.70 GPU-hours for DistilBERT, Llama~2, and Mistral, respectively. The probe totals include 0.64--0.67 GPU-hours of backbone feature extraction before head fitting. Thus, the main computational advantage of DistilBERT in this comparison is repeated inference rather than a large reduction in model-development cost. Appendix Table~A1d separates these stages and reports the repeated latency measurements.

The operational checkpoint's calibration diagnostics are reported next. Table~\ref{tab:calibration-quality} reports the completed DistilBERT results before and after temperature scaling. NLL evaluates the probability assigned to the true class and strongly penalizes confident mistakes; the Brier score evaluates the full three-class probability vector; and ECE summarizes the gap between confidence and observed accuracy across confidence groups. On the calibration partition, NLL fell from 0.1411 to 0.1041, Brier score from 0.0559 to 0.0514, fixed-width ECE from 0.0244 to 0.0083, and adaptive ECE from 0.0244 to 0.0138.

\begin{table*}[t]
\centering
\caption{DistilBERT calibration quality before and after temperature scaling}
\label{tab:calibration-quality}
\footnotesize
\setlength{\tabcolsep}{4pt}
\renewcommand{\arraystretch}{1.18}
\begin{tabularx}{\textwidth}{@{}l c c c c c@{}}
\toprule
Confidence score & Temperature & NLL & Brier & ECE & Adaptive ECE \\
\midrule
Raw MSP & 1.0000 & 0.1411 & 0.0559 & 0.0244 & 0.0244 \\
Temperature-scaled MSP & 1.7706 & 0.1041 & 0.0514 & 0.0083 & 0.0138 \\
\bottomrule
\end{tabularx}
\end{table*}

Figure~\ref{fig:calibration-routing} connects the three calibration decisions that support the operating policy. Panel (a) shows how the observed-accuracy curve moves toward the ideal confidence--accuracy diagonal after scaling; the vertical segments visualize the bin-wise calibration gaps summarized by ECE. Panel (b) places NLL, Brier score, and ECE on a common relative scale while retaining their original values in the annotations. Panel (c) then shows the risk--coverage behavior used to characterize the fixed routing policy.

\begin{figure*}[t]
\centering
\includegraphics[width=\textwidth]{figures/figure_a_calibration_and_routing.pdf}
\caption{DistilBERT calibration and selective-routing diagnostics. (a) Reliability before and after temperature scaling; ECE is the sample-weighted mean absolute gap between confidence and observed accuracy across bins, not a single plotted coordinate. (b) Raw and scaled NLL, Brier score, and ECE shown relative to each raw value (100\%); lower values indicate better calibration. (c) Risk--coverage curve with the fixed $\tau^*=0.70$ operating point.}
\label{fig:calibration-routing}
\end{figure*}

Using the calibrated confidence scores, candidates from 0.70 to 1.00 are swept on the dedicated calibration split and evaluated with the risk-coverage rule described in the Methodology section. The primary DistilBERT operating policy uses $T^*=1.7706$, $\alpha=0.05$, and $\tau^*=0.70$. The calibration-side upper-risk bound at this operating point was 2.82\%, satisfying the prespecified feasibility condition. Because $\tau=0.70$ was the lowest feasible candidate in the prespecified grid, it retained the greatest Phase~1 coverage; it was not chosen by examining held-out or Phase~2 outcomes.

Table~\ref{tab:routing-policy} connects the calibration-selected operating point to its held-out routing behavior. The final original-text Phase~2 experiment uses this prespecified policy at $\tau^*=0.70$; no threshold is chosen or revised using Phase~2 accuracy.

\begin{table*}[t]
\centering
\caption{Calibration-selected primary routing policy and held-out error concentration}
\label{tab:routing-policy}
\scriptsize
\setlength{\tabcolsep}{2.2pt}
\renewcommand{\arraystretch}{1.18}
\begin{tabularx}{\textwidth}{@{}c c c c c c c c c@{}}
\toprule
Risk $\alpha$ & $\tau^*$ & Cal. upper risk & Routed $n$ & Coverage & Routed P1 errors & Error capture & Routed P1 acc. & Error enrichment \\
\midrule
5.0\% & 0.70 & 2.82\% & 171 & 98.58\% & 87 & 21.91\% & 49.12\% & 15.38$\times$ \\
\bottomrule
\end{tabularx}
\vspace{2pt}
\parbox{\textwidth}{\footnotesize \textit{Note.} The risk budget and candidate grid were fixed before held-out and Phase~2 evaluation. The threshold was selected only on the 12,000-example calibration split. The fixed policy concentrated 87 of 397 Phase~1 errors in 171 routed posts. Error enrichment is the routed error rate (50.88\%) divided by the full-test error rate (3.31\%).}
\end{table*}

Routing effectiveness was assessed independently of the subsequent LLM outcomes. Under the primary policy, 171 posts (1.42\%) were routed. The routed subset contained 87 of 397 Phase~1 errors, giving 21.91\% error capture and 50.88\% routing precision. The corresponding full-test error rate was only 3.31\%, so the routed subset was approximately 15.4 times as error-dense as the full held-out set. Equivalently, a random subset of 171 posts would contain about 5.7 Phase~1 errors in expectation at the full-test error rate, whereas confidence routing concentrated 87 observed errors. Its 49.12\% Phase~1 accuracy was therefore far below the 96.69\% full-test accuracy, demonstrating that calibrated confidence concentrated an error-prone subset while sending only a small fraction of posts to Phase~2. Appendix Table~A4g reports the same decomposition for both evaluation sets.

Figure~\ref{fig:routing-concentration} visualizes this concentration mechanism on both evaluation sets. The confidence distributions explain why the same fixed threshold routes different proportions of Reddit and Mixed Emotion inputs. The error-capture curve compares confidence routing with a random-routing reference, while the final panel reports how much more error-dense each routed subset is than its complete evaluation set.

\begin{figure*}[t]
\centering
\includegraphics[width=\textwidth]{figures/figure_d_routing_concentration.pdf}
\caption{Confidence-based error concentration under the fixed routing policy. (a) Distribution of calibrated Phase~1 confidence with the same $\tau^*=0.70$ threshold applied to both evaluation sets. (b) Fraction of Phase~1 errors captured as increasingly large low-confidence subsets are routed; the dashed diagonal denotes random routing. (c) Overall and routed-subset Phase~1 error rates, with error-enrichment ratios shown above the routed values.}
\label{fig:routing-concentration}
\end{figure*}

\subsection{High-Confidence Accepted-Error Audit}

The complementary accepted-set audit examined the 11,829 predictions retained by Phase~1. Of these, 310 were errors, corresponding to 2.62\% accepted selective risk and 97.38\% accepted accuracy. Among 3,939 accepted reference-Depression posts, 125 were predicted as Happy or Neutral, yielding an accepted Depression false-negative risk of 3.17\%. Error frequency declined sharply with confidence, but did not vanish: 34 accepted errors, including 16 Depression false negatives, had calibrated confidence of at least 0.98.

\begin{table*}[t]
\centering
\caption{High-confidence accepted-error audit on the Reddit held-out test set}
\label{tab:accepted-error-audit}
\footnotesize
\setlength{\tabcolsep}{4pt}
\renewcommand{\arraystretch}{1.16}
\begin{tabularx}{\textwidth}{@{}l c c c c c c@{}}
\toprule
Reference proxy label & Accepted $n$ & Accepted errors & Accepted risk & To Depression & To Neutral & To Happy \\
\midrule
Depression & 3,939 & 125 & 3.17\% & -- & 10 & 115 \\
Neutral & 3,978 & 41 & 1.03\% & 17 & -- & 24 \\
Happy & 3,912 & 144 & 3.68\% & 112 & 32 & -- \\
\midrule
All classes & 11,829 & 310 & 2.62\% & 129 & 42 & 139 \\
\bottomrule
\end{tabularx}
\vspace{2pt}
\parbox{\textwidth}{\footnotesize \textit{Note.} ``Accepted risk'' is calculated within each accepted reference class. The last three columns count erroneous predictions only; correct predictions are omitted. The Depression false-negative risk is $125/3{,}939=3.17\%$. Reference labels are subreddit-derived proxies, so disagreements may reflect model errors or proxy-label/content mismatches.}
\end{table*}

Confidence-band analysis further localized the residual risk. Accepted error rates were 36.43\%, 25.95\%, 13.31\%, 4.18\%, and 0.36\% in the 0.70--0.80, 0.80--0.90, 0.90--0.95, 0.95--0.98, and 0.98--1.00 bands, respectively. Thus, confidence remained strongly informative, but no thresholded accepted set was error-free. Original-post review identified trajectory reversal and sarcasm, acute-distress cues overridden by a short positive clause, technical context masking explicit pride or excitement, and topic-term shortcuts triggered by mental-health vocabulary. Appendix Table~A4e summarizes these candidate error patterns, and Appendix Table~A4f supplies an original-post excerpt and the complete decision record for each selected case. The purpose of this audit is residual-risk transparency and failure analysis, not post hoc relabeling or evidence that Phase~2 produced an explanation; these examples were accepted by Phase~1 and therefore were never routed.

\subsection{Two-Phase System Efficacy and Confidence-Guided Routing}

With the classifier and routing policy fixed, RQ3 evaluates whether the assigned Phase~2 reasoner improves the complete system rather than merely producing plausible explanations on selected cases. Llama~2 CoT and Llama~3 SELF-DISCOVER therefore receive the same routed IDs within each evaluation set, and every corrected error is considered together with any newly introduced error. The following Reddit and Mixed Emotion subsections report the routed-only behavior, recombined full-set metrics, and paired statistical results without treating low confidence as a guarantee of successful correction.

\subsection{Reddit Held-Out Test: End-to-End Results}

Table~\ref{tab:reddit-e2e-results} reports the end-to-end comparison on the same 12,000-example Reddit held-out test set used for the routing analysis. Accepted examples retain the Phase~1 DistilBERT label, and only the 171 routed examples are eligible for Phase~2 replacement. The two reasoners receive the same minimally sanitized original title--body input. Within the routed subset, accuracy changed from 49.12\% under Phase~1 to 47.37\% with Llama~2 and 66.67\% with Llama~3. Accordingly, Llama~2 CoT produced a negligible $-0.03$-percentage-point full-test change (47 corrected versus 50 introduced routed errors), whereas Llama~3 SELF-DISCOVER increased full-test accuracy by 0.25 percentage points through 42 corrected and 12 introduced routed errors, or 30 net corrections. This increase is modest in full-test percentage points because only 1.42\% of the test set is routed, but the routed-only change is larger and the paired full-test effect is statistically supported.

\begin{table*}[t]
\centering
\caption{Reddit held-out end-to-end comparison under the final model-specific prompt policy}
\label{tab:reddit-e2e-results}
\scriptsize
\setlength{\tabcolsep}{2.1pt}
\renewcommand{\arraystretch}{1.16}
\begin{tabularx}{\textwidth}{@{}>{\raggedright\arraybackslash}X c c c c c c c@{}}
\toprule
System & Test $n$ & Routed $n$ & Phase~1 acc. & Routed-only acc. & End-to-end acc. & Change & Net corrections \\
\midrule
DistilBERT Phase~1 only & 12,000 & 171 & 96.69\% & 49.12\% & 96.69\% & -- & 0 \\
DistilBERT $\rightarrow$ Llama~2 CoT & 12,000 & 171 & 96.69\% & 47.37\% & 96.67\% & $-0.03$ pp & $-3$ \\
DistilBERT $\rightarrow$ Llama~3 SELF-DISCOVER & 12,000 & 171 & 96.69\% & 66.67\% & 96.94\% & $+0.25$ pp & $+30$ \\
\bottomrule
\end{tabularx}
\vspace{2pt}
\parbox{\textwidth}{\footnotesize \textit{Note.} Llama~2 CoT corrected 47 routed Phase~1 errors and introduced 50; Llama~3 SELF-DISCOVER corrected 42 and introduced 12. Both received the same minimally sanitized original text. Changes are calculated from exact counts before rounding, not by subtracting displayed accuracies. All metrics use the same 12,000 test examples.}
\end{table*}

\subsection{Handling Mixed Emotion Inputs: Quantitative Results}

To further examine model behavior under emotionally complex conditions, we evaluated the completed DistilBERT Phase~1 model on the 300-example supplementary synthetic Mixed Emotion Dataset. This dataset is not a primary benchmark or evidence of general clinical validity. Instead, it is a targeted stress test for cases in which positive, neutral, and depression-related cues co-occur or shift across the text. The Reddit-calibrated temperature and routing threshold were applied unchanged; neither was fitted or re-selected on the synthetic data.

DistilBERT achieved 81.33\% Phase~1 accuracy (244/300) on the final stress test. The model correctly classified 66 Depression examples, 95 Happy examples, and 83 Neutral examples. Of the 44 routed examples, 21 were Phase~1 errors, corresponding to 52.27\% routed-subset accuracy and 37.50\% capture of the 56 total Phase~1 errors. The routed error rate was 47.73\%, compared with 18.67\% across all 300 examples, a 2.56-fold concentration. This behavior is consistent with the intended role of the set: it concentrates evaluation on emotionally mixed and trajectory-sensitive examples rather than estimating prevalence-level performance.

Llama~2 CoT increased end-to-end accuracy by 4.00 percentage points (81.33\% to 85.33\%), correcting 18 of the 21 routed Phase~1 errors while introducing six. Llama~3 SELF-DISCOVER increased accuracy by 6.00 percentage points (81.33\% to 87.33\%), also correcting 18 errors (85.71\%) but introducing none. Routed-only accuracy was 79.55\% and 93.18\%, respectively, versus 52.27\% for Phase~1. Both configurations used the same 300 inputs, fixed Reddit-calibrated policy, and 44 routed cases.

\begin{table*}[t]
\centering
\caption{Mixed Emotion stress-test end-to-end results on the shared routed subset}
\label{tab:mixed-emotion-results}
\scriptsize
\setlength{\tabcolsep}{2.1pt}
\renewcommand{\arraystretch}{1.16}
\begin{tabularx}{\textwidth}{@{}>{\raggedright\arraybackslash}X c c c c c c c@{}}
\toprule
System & Routed $n$ & Routed-only acc. & Corrected & Introduced & Net & End-to-end acc. & Change \\
\midrule
DistilBERT Phase~1 only & 44 & 52.27\% & -- & -- & -- & 81.33\% & -- \\
DistilBERT $\rightarrow$ Llama~2 CoT & 44 & 79.55\% & 18 & 6 & 12 & 85.33\% & $+4.00$ pp \\
DistilBERT $\rightarrow$ Llama~3 SELF-DISCOVER & 44 & 93.18\% & 18 & 0 & 18 & 87.33\% & $+6.00$ pp \\
\bottomrule
\end{tabularx}
\vspace{2pt}
\parbox{\textwidth}{\footnotesize \textit{Note.} All systems use the same 300 supplementary inputs and the same 44 routed examples. Phase~1 accuracy is 81.33\%. ``Corrected'' and ``introduced'' count changes only within the routed subset; ``net'' is corrected minus introduced.}
\end{table*}

Figure~\ref{fig:end-to-end-effect} summarizes the distinction between identifying difficult inputs and correcting them. Panel (a) reports accuracy only on the routed subsets, panel (b) shows the resulting change after routed predictions are recombined with accepted Phase~1 predictions, and panel (c) separates corrected from introduced errors. This decomposition prevents a positive correction count from obscuring errors newly introduced by a reasoner.

\begin{figure*}[t]
\centering
\includegraphics[width=\textwidth]{figures/figure_b_end_to_end_effect.pdf}
\caption{Effect of selective LLM re-evaluation on the shared routed subsets and complete evaluation sets. (a) Phase~1 and Phase~2 accuracy on routed examples. (b) Full-set accuracy change after recombination with accepted Phase~1 predictions. (c) Corrected and introduced routed errors; net improvement depends on their difference rather than on corrections alone.}
\label{fig:end-to-end-effect}
\end{figure*}

Table~\ref{tab:paired-statistics} reports paired uncertainty analyses for the same end-to-end comparisons. On Reddit, the Llama~2 confidence interval includes zero and its exact McNemar test does not distinguish the end-to-end result from the Phase~1 baseline. In contrast, the original-text Llama~3 interval is entirely positive, and its improvement remains significant after Holm adjustment across the four comparisons. On the Mixed Emotion stress test, both paired bootstrap intervals are positive and both changes remain significant after Holm adjustment; Llama~3 again provides the larger net improvement and introduces no routed error.

\begin{table*}[t]
\centering
\caption{Paired statistical analysis of Phase~1 and end-to-end correctness}
\label{tab:paired-statistics}
\footnotesize
\setlength{\tabcolsep}{3pt}
\renewcommand{\arraystretch}{1.16}
\begin{tabularx}{\textwidth}{@{}l >{\raggedright\arraybackslash}X c c c c c c@{}}
\toprule
Dataset & Reasoner & Change & Corrected & Introduced & Paired 95\% CI & Exact $p$ & Holm $p$ \\
\midrule
Reddit & Llama~2 CoT & $-0.03$ pp & 47 & 50 & [$-0.18$, 0.13] & 0.839 & 0.839 \\
Reddit & Llama~3 SELF-DISCOVER & $+0.25$ pp & 42 & 12 & [0.13, 0.38] & $<0.0001$ & 0.0002 \\
Mixed Emotion & Llama~2 CoT & $+4.00$ pp & 18 & 6 & [1.00, 7.33] & 0.0227 & 0.0453 \\
Mixed Emotion & Llama~3 SELF-DISCOVER & $+6.00$ pp & 18 & 0 & [3.33, 8.67] & $<0.0001$ & $<0.0001$ \\
\bottomrule
\end{tabularx}
\vspace{2pt}
\parbox{\textwidth}{\footnotesize \textit{Note.} Confidence intervals are paired percentile-bootstrap intervals from 50,000 resamples. Exact $p$-values are two-sided McNemar tests over corrected and introduced pairs. Holm $p$ adjusts the four comparisons shown. The Reddit analysis follows the manuscript's exact 12,000-row balanced protocol.}
\end{table*}

On Mixed Emotion, class-wise Llama~3 recall increased from 66\% to 79\% for Depression (13 net corrections), 95\% to 99\% for Happy (four), and 83\% to 84\% for Neutral (one). Each class has 100 examples. No class incurred an introduced error in this run. Appendix Table~A4c separates clear informational scenarios from mixed-cue and trajectory-shift scenarios.

Figure~\ref{fig:error-structure} shows the corresponding class-level error structure. The Reddit matrices reveal reductions in both Depression-to-Happy and Happy-to-Depression confusion after Llama~3 re-evaluation. The Mixed Emotion matrices show that the largest change is recovery of Depression examples initially assigned to Happy, while the Neutral-to-Happy residual remains unchanged under the fixed router.

\begin{figure*}[t]
\centering
\includegraphics[width=0.86\textwidth]{figures/figure_c_error_structure.pdf}
\caption{Confusion matrices before and after selective Llama~3 re-evaluation. Cells report counts and row-normalized percentages. (a)--(b) Reddit held-out test. (c)--(d) Mixed Emotion stress test. Accepted predictions remain unchanged; differences arise only from the fixed routed subsets.}
\label{fig:error-structure}
\end{figure*}

\subsection{Additional Same-Source Holdout Results}
\label{sec:additional-holdout-results}

On the additional 9,000 posts, the unchanged operational classifier achieved 96.73\% accuracy and routed 137 posts (1.52\%). The routed subset contained 67 of 294 Phase~1 errors (22.79\% error capture), with a 48.91\% error rate versus 3.27\% overall. The remaining 8,863 accepted posts had a 2.56\% error rate. Thus, the fixed policy again concentrated errors in a small subset, while leaving 227 accepted errors outside the correction path.

Table~\ref{tab:additional-holdout} reports the complete-set effects. Llama~3 SELF-DISCOVER corrected 36 errors and introduced nine, giving 27 net corrections and 97.03\% accuracy. Its $+0.30$ percentage-point gain remained significant after Holm adjustment ($p=0.000131$). Llama~2 CoT corrected 41 errors but introduced 26, yielding 96.90\% accuracy; its paired interval included zero. This distinction supports the evaluated Llama~3 configuration without implying that any LLM re-evaluator will improve the classifier or establishing a direct statistical difference between the two reasoners.

\begin{table*}[t]
\centering
\caption{Additional same-source holdout: fixed-policy performance on 9,000 posts}
\label{tab:additional-holdout}
\footnotesize
\setlength{\tabcolsep}{4pt}
\renewcommand{\arraystretch}{1.16}
\begin{tabularx}{\textwidth}{@{}Y c c c c c c c@{}}
\toprule
Configuration & Accuracy & Macro F1 & Corrected & Introduced & Change & 95\% CI (pp) & Holm $p$ \\
\midrule
Phase~1 & 96.73\% & 96.73\% & -- & -- & -- & -- & -- \\
+ Llama~2 CoT & 96.90\% & 96.90\% & 41 & 26 & $+0.17$ pp & [$-0.01$, 0.34] & 0.0864 \\
+ Llama~3 SELF-DISCOVER & 97.03\% & 97.03\% & 36 & 9 & $+0.30$ pp & [0.16, 0.46] & 0.000131 \\
\bottomrule
\end{tabularx}
\vspace{2pt}
\parbox{\textwidth}{\footnotesize\textit{Note.} Both reasoners re-evaluated the same 137 posts; all other predictions were retained. No parsing failures occurred. Changes and paired percentile-bootstrap intervals use all 9,000 posts (10,000 resamples). Holm adjustment covers the two exact McNemar comparisons against Phase~1 in this additional evaluation, separately from Table~\ref{tab:paired-statistics}.}
\end{table*}

The Llama~3 improvement corresponds to a 9.18\% reduction in total errors (294 to 267). This evaluation extends the evidence beyond the prompt-development sample, while remaining conditional on the same source and proxy-label construction. Appendix~F reports class-wise results; improvement was not uniform across all classes.

\subsection{Observed Correction Cases and Reasoning Traceability}

Appendix~D reports one observed Mixed Emotion Depression correction for each Phase~2 model using the stored output columns rather than an invented illustration. It also reports an observed Reddit original-text Llama~3 correction in which Phase~1 predicted Happy for a post describing numbness and exhaustion from trying to remain positive. Llama~2 CoT retains independent assessment, Phase~1 comparison, and terminal-label fields, while Llama~3 SELF-DISCOVER retains an auditable SELECT--ADAPT--IMPLEMENT trace before the final answer.

The examples show how generated responses relate local cues to the complete post. Only the parsed terminal label enters quantitative evaluation; the rationale remains available for inspection. Corrected and introduced counts establish the observed predictive effect, whereas the trace records a generated justification rather than a verified causal account of the decision.

\subsection{Interpretation of Routing and Re-Evaluation Results}

The three research questions receive distinct answers. For \textbf{RQ1a}, DistilBERT achieved strong held-out predictive performance, and temperature scaling reduced NLL, Brier score, and fixed-bin ECE on the calibration split. For \textbf{RQ1b}, the completed three-seed comparison retained DistilBERT: its mean held-out accuracy and macro F1 were both 96.88\%, compared with 95.36\% for the Mistral frozen-backbone probe and 95.17\% for the Llama~2 probe. This is evidence for the operational choice under the bounded comparison, reinforced by the measured inference advantage, not a ranking of maximally tuned architectures. For \textbf{RQ2}, the fixed confidence policy preserved 98.58\% Reddit coverage while concentrating 87 of the 397 Phase~1 errors in the 171 routed posts; routed-sample accuracy was 49.12\%, compared with 96.69\% overall. The same fixed policy identified a 44-example Mixed Emotion subset with 52.27\% Phase~1 accuracy, compared with 81.33\% overall. For \textbf{RQ3}, re-evaluation benefit was model-dependent. Llama~2 remained statistically indistinguishable from Phase~1 on Reddit, whereas Llama~3 raised routed-only accuracy to 66.67\% and produced a positive paired full-test effect. On the controlled stress test, Llama~2 and Llama~3 raised routed-only accuracy to 79.55\% and 93.18\%, respectively, with Llama~3 avoiding the six introduced errors observed for Llama~2. Confidence therefore identifies a correction opportunity, but the reasoning policy determines how much of that opportunity is realized.

Only 1.42\% of the original Reddit test posts, 14.67\% of Mixed Emotion examples, and 1.52\% of additional holdout posts underwent re-evaluation. These are fractions of inputs, not counts of individual LLM calls: CoT and SELF-DISCOVER require multiple generation stages per routed post. The Phase~1 benchmark supports the classifier's inference advantage, but the fraction avoiding re-evaluation is not a measured GPU-time saving because every input incurs Phase~1 computation and no all-input LLM timing baseline was run. The stress-test routing rate characterizes its constructed scenario mix, not an expected deployment workload.

A conditional routing oracle further clarifies the scale of the observed gains without introducing a new tuned baseline. If every routed Phase~1 error were corrected and no correct routed prediction were changed, Reddit accuracy could rise from 96.69\% to at most 97.42\% under the fixed 171-post routing policy; the 30 Llama~3 net corrections realize 34.5\% of the 87-error correction opportunity. On Mixed Emotion, the corresponding ceiling is 88.33\%, and the 18 Llama~3 net corrections realize 85.7\% of the 21-error opportunity. These ceilings are conditional on the already selected router and are not claims about an attainable model. Their purpose is to distinguish a limited routing opportunity from a failure to exploit routed cases. Definitions and the complete decomposition are reported in Appendix Table~A4g.

In addition to selective computation, the framework adds an inspectable rationale channel for routed cases. Unlike a single-stage classifier that outputs only a class label or probability score, the Phase 2 record retains the generated rationale and terminal label. This improves case-level traceability, but the rationale is not treated as proof of faithfulness, clinical validity, or prediction correctness.

Taken together, the experiments support the Phase~1 and routing components of the framework for the constructed proxy-label task. Llama~3 produced positive paired effects on the original Reddit test, the Mixed Emotion stress test, and the additional same-source holdout. Llama~2 showed no statistically detectable gain on either Reddit evaluation. The additional holdout strengthens evidence for the frozen Llama~3 configuration beyond the prompt-development sample, but does not establish broad real-world mental-health generalization or isolate the effect of the reasoning method from the underlying model.

\section{Limitations and Future Work}

\textit{Label and population validity.} Reddit labels reflect subreddit membership and polarity filtering, not clinical assessments. Filtering may favor sentiment-aligned posts and make the task easier than classification of unfiltered discourse. The balanced, post-level split also does not establish performance on unseen authors, future time periods, or other platforms. The synthetic stress test has controlled labels and scenarios but may contain generator-specific phrasing and repeated narrative patterns. Neither source establishes clinical validity. Independent annotation is needed to assess post-level label consistency and rationale quality before considering any decision-support use.

\textit{Residual risk and calibration.} The router left 310 errors among 11,829 accepted predictions, including 34 at confidence 0.98 or above. Confidence therefore supports selection without certifying individual predictions. Temperature and threshold selection share a calibration partition, and the risk bound is a pointwise selection criterion rather than a post-selection guarantee. Results under the fixed Reddit policy do not establish risk control on other distributions. The accepted-error examples identify candidate failure patterns; without attribution tests or independent adjudication, they do not establish causal mechanisms or failure-mode prevalence.

\textit{Comparison scope.} The Phase~1 study compares three practical configurations under a limited search budget, not equivalently adapted or maximally optimized architectures. Three seeds characterize a narrow range of training variability. In Phase~2, model family and prompting method vary together, so the observed advantage of Llama~3 SELF-DISCOVER cannot be attributed to SELF-DISCOVER alone. A crossed model--prompt evaluation would be needed to separate these effects. Similarly, the final original-text rerun supports the reported input policy but does not establish a general causal effect of preserving any particular textual feature.

\textit{Prompt development and statistical scope.} Prompt variants were examined on the original routed Reddit cases under the earlier cleaned-input condition. Freezing prompts before the original-text rerun did not remove that development dependence, so those original end-to-end comparisons remain exploratory. The additional 9,000-post holdout evaluates the frozen pipeline on previously unused same-source text and provides separate evidence of improvement, without retroactively changing the status of the original analyses. Normalized exact-text exclusion does not establish author independence or remove all near-duplicates, and the paired intervals do not quantify repeated-generation variability. External or author-disjoint evaluation and expert-reviewed labels remain important extensions. This issue is distinct from calibration-partition fitting of $T^*$ and $\tau^*$.

\textit{Computational and explanatory scope.} The controlled Phase~1 benchmark compares our implementations on one accelerator model and serving configuration, not optimized inference systems across hardware platforms. Three timing repetitions characterize within-session variation; they do not capture all variation across Colab sessions. GPU-board energy estimates exclude CPU and system-level consumption and have limited temporal resolution for short timed regions. Total cascade energy and monetary savings were not established. Our per-post SELF-DISCOVER adaptation does not reuse a task-level plan, so efficiency results from the original method cannot be transferred to this implementation. Generated rationales still require independent assessment of faithfulness and label support.

\section{Conclusion}

This study evaluated selective LLM re-evaluation by separating classifier performance, calibration, error concentration, and net correction. The three-seed comparison and common-hardware inference benchmark supported retaining DistilBERT within the tested configurations: it combined the highest mean predictive scores with substantially higher batched throughput, lower single-post latency, and lower inference memory. Its fixed operational checkpoint routed 1.42\% of Reddit posts, capturing 87 of 397 Phase~1 errors. Llama~3 SELF-DISCOVER corrected 42 errors and introduced 12, raising full-set accuracy from 96.69\% to 96.94\%; Llama~2 CoT produced no detectable improvement. On the synthetic stress test, the corresponding gains were 6.00 and 4.00 percentage points.

The central finding is that selecting an error-prone subset and correcting it are separate requirements. On an additional 9,000-post same-source holdout, the frozen pipeline re-evaluated 1.52\% of posts and improved accuracy from 96.73\% to 97.03\% with Llama~3, with a significant paired gain. This strengthens support beyond the original prompt-development sample while leaving accepted errors outside the correction path. External validation and expert-reviewed labels are needed to establish broader applicability. The framework remains a proxy emotion-classification pipeline, not a clinical diagnostic system.

\balance
\begin{thebibliography}{46}

\bibitem{ref1} World Health Organization, ``Depressive disorder (depression),'' WHO Fact Sheet, Aug. 29, 2025. Available: \url{https://www.who.int/news-room/fact-sheets/detail/depression}

\bibitem{ref2} World Health Organization, ``Depression: let's talk' says WHO, as depression tops list of causes of ill health,'' WHO News Release, Mar. 30, 2017. Available: \url{https://www.who.int/news/item/30-03-2017--depression-let-s-talk-says-who-as-depression-tops-list-of-causes-of-ill-health}

\bibitem{ref3} World Health Organization, ``Mental health at work,'' WHO Fact Sheet, Sept. 2, 2024. Available: \url{https://www.who.int/news-room/fact-sheets/detail/mental-health-at-work}

\bibitem{ref4} World Health Organization, ``COVID-19 pandemic triggers 25\% increase in prevalence of anxiety and depression worldwide,'' WHO News Release, Mar. 2, 2022. Available: \url{https://www.who.int/news/item/02-03-2022-covid-19-pandemic-triggers-25-increase-in-prevalence-of-anxiety-and-depression-worldwide}

\bibitem{ref5} B. Levis, A. Benedetti, B. D. Thombs, and the DEPRESsion Screening Data (DEPRESSD) Collaboration, ``Accuracy of Patient Health Questionnaire-9 (PHQ-9) for screening to detect major depression: Individual participant data meta-analysis,'' BMJ, vol. 365, l1476, 2019. doi: 10.1136/bmj.l1476

\bibitem{ref6} OECD, A New Benchmark for Mental Health Systems: Tackling the Social and Economic Costs of Mental Ill-Health. OECD Health Policy Studies, OECD Publishing, Paris, 2021. doi: 10.1787/4ed890f6-en

\bibitem{ref7} R. A. Calvo, D. N. Milne, M. S. Hussain, and H. Christensen, ``Natural language processing in mental health applications using non-clinical texts,'' Natural Language Engineering, vol. 23, no. 5, pp. 649-685, 2017. doi: 10.1017/S1351324916000383

\bibitem{ref8} S. Ji, T. Zhang, L. Ansari, J. Fu, P. Tiwari, and E. Cambria, ``MentalBERT: Publicly available pretrained language models for mental healthcare,'' in Proceedings of the Thirteenth Language Resources and Evaluation Conference, Marseille, France, pp. 7184-7190, European Language Resources Association, 2022. Available: \url{https://aclanthology.org/2022.lrec-1.778/}

\bibitem{ref9} T. Zhang, K. Yang, S. Ji, and S. Ananiadou, ``Emotion fusion for mental illness detection from social media: A survey,'' Information Fusion, vol. 92, pp. 231-246, 2023. doi: 10.1016/j.inffus.2022.11.031

\bibitem{ref10} R. Kumar, K. Maharaj, A. Saxena, and P. Bhattacharyya, ``Mental disorder classification via temporal representation of text,'' in Findings of the Association for Computational Linguistics: EMNLP 2024, Miami, FL, USA, pp. 10901-10916, Association for Computational Linguistics, 2024. doi: 10.18653/v1/2024.findings-emnlp.639

\bibitem{ref11} S. Tonekaboni, S. Joshi, M. D. McCradden, and A. Goldenberg, ``What clinicians want: Contextualizing explainable machine learning for clinical end use,'' in Proceedings of the 4th Machine Learning for Healthcare Conference, PMLR, vol. 106, pp. 359-380, 2019. Available: \url{https://proceedings.mlr.press/v106/tonekaboni19a.html}

\bibitem{ref17} J. Wei et al., ``Chain-of-thought prompting elicits reasoning in large language models,'' in Advances in Neural Information Processing Systems, vol. 35, pp. 24824-24837, 2022. doi: 10.48550/arXiv.2201.11903

\bibitem{ref18} P. Zhou et al., ``SELF-DISCOVER: Large language models self-compose reasoning structures,'' in Advances in Neural Information Processing Systems, vol. 37, 2024. doi: 10.52202/079017-4004

\bibitem{ref20} S. Rude, E. M. Gortner, and J. Pennebaker, ``Language use of depressed and depression-vulnerable college students,'' Cognition and Emotion, vol. 18, no. 8, pp. 1121-1133, 2004. doi: 10.1080/02699930441000030

\bibitem{ref21} M. De Choudhury, M. Gamon, S. Counts, and E. Horvitz, ``Predicting depression via social media,'' Proceedings of the International AAAI Conference on Web and Social Media, vol. 7, no. 1, pp. 128-137, 2013. doi: 10.1609/icwsm.v7i1.14432

\bibitem{ref22} S. Tsugawa, Y. Kikuchi, F. Kishino, K. Nakajima, Y. Itoh, and H. Ohsaki, ``Recognizing depression from Twitter activity,'' in Proceedings of the 33rd Annual ACM Conference on Human Factors in Computing Systems (CHI), ACM, pp. 3187-3196, 2015. doi: 10.1145/2702123.2702280

\bibitem{ref23} G. Coppersmith, M. Dredze, C. Harman, K. Hollingshead, and M. Mitchell, ``CLPsych 2015 shared task: Depression and PTSD on Twitter,'' in Proceedings of the 2nd Workshop on Computational Linguistics and Clinical Psychology: From Linguistic Signal to Clinical Reality, Denver, CO, pp. 31-39, Association for Computational Linguistics, 2015. doi: 10.3115/v1/W15-1204

\bibitem{ref29} G. Shen, J. Jia, L. Nie, F. Feng, C. Zhang, T. Hu, T.-S. Chua, and W. Zhu, ``Depression detection via harvesting social media: A multimodal dictionary learning solution,'' in Proceedings of the 26th International Joint Conference on Artificial Intelligence (IJCAI), Melbourne, Australia, pp. 3838-3844, 2017. doi: 10.24963/ijcai.2017/536

\bibitem{ref31} M. M. Tadesse, H. Lin, B. Xu, and L. Yang, ``Detection of depression-related posts in Reddit social media forum,'' IEEE Access, vol. 7, pp. 44883-44893, 2019. doi: 10.1109/ACCESS.2019.2909180

\bibitem{ref24} A. Wongkoblap, M. Vadillo, and V. Curcin, ``Researching mental health disorders in the era of social media: Systematic review,'' Journal of Medical Internet Research, vol. 19, no. 6, e228, 2017. doi: 10.2196/jmir.7215

\bibitem{ref25} T. Mikolov, K. Chen, G. Corrado, and J. Dean, ``Efficient estimation of word representations in vector space,'' Proceedings of Workshop at ICLR, arXiv preprint arXiv:1301.3781, 2013. doi: 10.48550/arXiv.1301.3781

\bibitem{ref26} J. Pennington, R. Socher, and C. D. Manning, ``GloVe: Global vectors for word representation,'' in Proceedings of the 2014 Conference on Empirical Methods in Natural Language Processing (EMNLP), pp. 1532-1543, 2014. doi: 10.3115/v1/D14-1162

\bibitem{ref27} S. Hochreiter and J. Schmidhuber, ``Long short-term memory,'' Neural Computation, vol. 9, no. 8, pp. 1735-1780, 1997. doi: 10.1162/neco.1997.9.8.1735

\bibitem{ref28} Y. Kim, ``Convolutional neural networks for sentence classification,'' in Proceedings of the 2014 Conference on Empirical Methods in Natural Language Processing (EMNLP), Doha, Qatar, pp. 1746-1751, Association for Computational Linguistics, 2014. doi: 10.3115/v1/D14-1181

\bibitem{ref30} A. H. Orabi, P. Buddhitha, M. H. Orabi, and D. Inkpen, ``Deep learning for depression detection of Twitter users,'' in Proceedings of the 5th Workshop on Computational Linguistics and Clinical Psychology (CLPsych), Association for Computational Linguistics, pp. 88-97, 2018. doi: 10.18653/v1/W18-0609

\bibitem{ref32} D. E. Losada, F. Crestani, and J. Parapar, ``Overview of eRisk 2018: Early risk prediction on the internet,'' CLEF 2018 Working Notes, CEUR Workshop Proceedings, vol. 2125, 2018. Available: \url{https://ceur-ws.org/Vol-2125/invited_paper_1.pdf}

\bibitem{ref33} Z. Zhang, J. Zhu, Z. Guo, Y. Zhang, Z. Li, and B. Hu, ``Natural language processing for depression prediction on Sina Weibo: Method study and analysis,'' JMIR Mental Health, vol. 11, e58259, 2024. doi: 10.2196/58259

\bibitem{ref34} I. Tavchioski, M. Robnik-\v{S}ikonja, and S. Pollak, ``Detection of depression on social networks using transformers and ensembles,'' arXiv preprint arXiv:2305.05325, 2023. doi: 10.48550/arXiv.2305.05325

\bibitem{ref35} D. Shin, H. Kim, S. Lee, Y. Cho, and W. Jung, ``Using large language models to detect depression from user-generated diary text data as a novel approach in digital mental health screening: Instrument validation study,'' Journal of Medical Internet Research, vol. 26, e54617, 2024. doi: 10.2196/54617

\bibitem{ref43} Y. Wang, D. Inkpen, and P. Kirinde Gamaarachchige, ``Explainable depression detection using large language models on social media data,'' in Proceedings of the 9th Workshop on Computational Linguistics and Clinical Psychology (CLPsych 2024), St. Julians, Malta, pp. 108-126, Association for Computational Linguistics, 2024. doi: 10.18653/v1/2024.clpsych-1.8

\bibitem{ref44} X. Lan, Z. Han, Y. Cheng, L. Sheng, J. Feng, C. Gao, and Y. Li, ``Depression detection on social media with large language models,'' in Proceedings of the 2025 Conference on Empirical Methods in Natural Language Processing: Industry Track, pp. 2155-2171, Association for Computational Linguistics, 2025. doi: 10.18653/v1/2025.emnlp-industry.151

\bibitem{ref45} M. Omar and I. Levkovich, ``Exploring the efficacy and potential of large language models for depression: A systematic review,'' Journal of Affective Disorders, vol. 371, pp. 234-244, 2025. doi: 10.1016/j.jad.2024.11.052

\bibitem{ref36} R. Salas-Z\'{a}rate, G. Alor-Hern\'{a}ndez, M. del P. Salas-Z\'{a}rate, M. A. Paredes-Valverde, M. Bustos-L\'{o}pez, and J. L. S\'{a}nchez-Cervantes, ``Detecting depression signs on social media: A systematic literature review,'' Healthcare (Basel), vol. 10, no. 2, 291, 2022. doi: 10.3390/healthcare10020291

\bibitem{ref37} S. Chancellor and M. De Choudhury, ``Methods in predictive techniques for mental health status on social media: A critical review,'' npj Digital Medicine, vol. 3, 43, 2020. doi: 10.1038/s41746-020-0233-7

\bibitem{ref46} Y. Cao, J. Dai, Z. Wang, Y. Zhang, X. Shen, Y. Liu, and Y. Tian, ``Machine learning approaches for mental illness detection on social media: A systematic review of biases and methodological challenges,'' Journal of Behavioral Data Science, vol. 5, no. 1, pp. 67-102, 2025. doi: 10.35566/jbds/caoyc

\bibitem{ref15} C. Guo, G. Pleiss, Y. Sun, and K. Q. Weinberger, ``On calibration of modern neural networks,'' in Proceedings of the 34th International Conference on Machine Learning, PMLR, vol. 70, pp. 1321-1330, 2017. Available: \url{https://proceedings.mlr.press/v70/guo17a.html}

\bibitem{ref16} Y. Geifman and R. El-Yaniv, ``Selective classification for deep neural networks,'' in Advances in Neural Information Processing Systems, vol. 30, pp. 4878-4887, 2017. Available: \url{https://proceedings.neurips.cc/paper/2017/hash/4a8423d5e91fda00bb7e46540e2b0cf1-Abstract.html}

\bibitem{ref38} M. Conway and D. O'Connor, ``Social media, big data, and mental health: Current advances and ethical implications,'' Current Opinion in Psychology, vol. 9, pp. 77-82, 2016. doi: 10.1016/j.copsyc.2016.01.004

\bibitem{ref39} A. Benton, G. Coppersmith, and M. Dredze, ``Ethical research protocols for social media health research,'' in Proceedings of the First ACL Workshop on Ethics in Natural Language Processing, Valencia, Spain, pp. 94-102, Association for Computational Linguistics, 2017. doi: 10.18653/v1/W17-1612

\bibitem{ref41} S. Loria, ``TextBlob documentation: Quickstart---sentiment analysis.'' Available: \url{https://textblob.readthedocs.io/en/dev/quickstart.html#sentiment-analysis}. Accessed: Sept. 9, 2026.

\bibitem{ref42} T. De Smedt and W. Daelemans, ``Pattern for Python,'' Journal of Machine Learning Research, vol. 13, no. 66, pp. 2063-2067, 2012.

\bibitem{ref19} A. Kang, J. Y. Chen, Z. Lee-Youngzie, and S. Fu, ``Synthetic data generation with LLM for improved depression prediction,'' arXiv preprint arXiv:2411.17672, 2024. doi: 10.48550/arXiv.2411.17672

\bibitem{ref12} V. Sanh, L. Debut, J. Chaumond, and T. Wolf, ``DistilBERT, a distilled version of BERT: Smaller, faster, cheaper and lighter,'' arXiv preprint arXiv:1910.01108, 2019. doi: 10.48550/arXiv.1910.01108

\bibitem{ref14} H. Touvron et al., ``Llama 2: Open foundation and fine-tuned chat models,'' arXiv preprint arXiv:2307.09288, 2023. doi: 10.48550/arXiv.2307.09288

\bibitem{ref40} A. Grattafiori et al., ``The Llama 3 Herd of Models,'' arXiv preprint arXiv:2407.21783, 2024. doi: 10.48550/arXiv.2407.21783

\bibitem{ref13} A. Q. Jiang et al., ``Mistral 7B,'' arXiv preprint arXiv:2310.06825, 2023. doi: 10.48550/arXiv.2310.06825
\end{thebibliography}

\clearpage
\onecolumn
\appendices
% Keep appendix headings and their supporting tables compact without orphaning titles.
\setlength{\medskipamount}{4pt plus 1pt minus 1pt}
\Needspace{12\baselineskip}
\section{Supplementary Methodological Details}

Appendix Table~A1 summarizes the main variables retained from the Reddit dataset during data construction and preprocessing. Appendix Table~A1b complements the schema description with actual class-balanced examples from the supplementary Mixed Emotion stress test.

\par\medskip
\noindent\textsc{Appendix Table A1. Description of Reddit dataset variables}\par
\label{tab:appendix-a1-reddit-variables}
\smallskip
\small
\setlength{\tabcolsep}{4pt}
\renewcommand{\arraystretch}{1.18}
\noindent\begin{tabularx}{\textwidth}{@{}>{\raggedright\arraybackslash}p{0.25\textwidth}Y@{}}
\toprule
Variable & Description \\
\midrule
Subreddit & The subreddit name where the post was published, indicating the context, topic, and implicit emotional milieu of the post. \\
Author & The anonymized Reddit username of the post's author. \\
\path{Over_18} & A binary flag indicating whether the post was tagged as adult content. \\
Title & The title of the Reddit post, often providing a brief summary of the content. \\
Selftext & The main body of the post, containing the full text written by the user, if available. \\
\path{Created_UTC} & The timestamp, in UTC, indicating when the post was created; this variable can support temporal analysis. \\
\path{Title_with_selftext} & The concatenation of the title and body text, representing the complete textual content of the post before cleaning. \\
\path{Title_with_selftext_cleaned} & The processed text after cleaning, used as the primary input for modeling. \\
Polarity & A document-level TextBlob polarity score computed by passing the cleaned concatenated title and body to \texttt{TextBlob(...).sentiment.polarity}; it is not a manually averaged token score. \\
\path{Class_group} & The target label denoting the emotional class: Depression, Neutral, or Happy. \\
\bottomrule
\end{tabularx}
\par\medskip

The examples below make the stress-test construct observable: the target is determined by the dominant full-text trajectory rather than by the presence of any single positive or negative word.

\Needspace{30\baselineskip}
\noindent\textsc{Appendix Table A1b. Representative inputs from the Mixed Emotion stress test}\par
\label{tab:appendix-a1b-mixed-examples}
\smallskip
\small
\setlength{\tabcolsep}{4pt}
\renewcommand{\arraystretch}{1.18}
\noindent\begin{tabularx}{\textwidth}{@{}>{\raggedright\arraybackslash}p{0.13\textwidth}>{\raggedright\arraybackslash}p{0.09\textwidth}>{\raggedright\arraybackslash}p{0.51\textwidth}Y@{}}
\toprule
Example ID & Target & Actual input text & Dominant-label rationale \\
\midrule
\texttt{MEV2\_DEP\_001} & Depression & During a school day, I noticed a kind message from a friend, and ordinary tasks still needed to be finished. Even with those ordinary or positive moments, there was a heavy sadness returning in the background. I kept rereading the same message without knowing how to answer, especially before breakfast. By the end, the positive parts felt brief, and the heavier sadness was still the feeling I carried with me. & Positive and routine details coexist with the negative cues, but the ending preserves unresolved sadness as the dominant trajectory. \\
\texttt{MEV2\_NEU\_001} & Neutral & While reviewing a class registration page, I noticed that one course appeared in two sections with different room numbers. I copied both entries into a note, checked the department page, and left the final choice blank until the schedule page is updated. The post is mainly about checking inconsistent information and does not state a strong personal reaction. & A minor procedural ambiguity is present, but the post remains factual and low in affective intensity. \\
\texttt{MEV2\_HAP\_001} & Happy & While dealing with a school day, I could still feel a heavy sadness returning in the background. At the same time, there was a kind message from a friend, and ordinary tasks still needed to be finished. I kept rereading the same message without knowing how to answer, especially before breakfast. Even with the mixed parts, I ended the day feeling relieved, quietly happy, and more connected than before. & The input contains sadness and routine details, but its explicit final resolution is relief, happiness, and renewed connection. \\
\bottomrule
\end{tabularx}
\par\smallskip
\normalsize

\noindent\footnotesize\textit{Note.} These are unmodified input texts from the final 300-example supplementary dataset. They were not used for Phase~1 training, temperature fitting, or threshold selection. The table illustrates co-occurring cues and final-trajectory labeling; it is not intended to establish population prevalence or clinical validity.\normalsize
\par\medskip

\Needspace{11\baselineskip}
\noindent\textsc{Appendix Table A1c. Sweep-selected Phase~1 configurations}\par
\label{tab:appendix-a1c-phase1-configurations}
\smallskip
\footnotesize
\setlength{\tabcolsep}{5pt}
\renewcommand{\arraystretch}{1.14}
\noindent\begin{tabularx}{\textwidth}{@{}>{\raggedright\arraybackslash}p{0.19\textwidth}>{\raggedright\arraybackslash}Y c c c c@{}}
\toprule
Model & Training & Learning rate & Batch & Epochs & Weight decay \\
\midrule
DistilBERT & Full fine-tuning & $4.037\times10^{-5}$ & 64 & 3 & $10^{-3}$ \\
Llama~2~7B & Linear probe & $4.151\times10^{-4}$ & 32 & 3 & $10^{-4}$ \\
Mistral~7B & Linear probe & $2.399\times10^{-4}$ & 64 & 2 & $10^{-4}$ \\
\bottomrule
\end{tabularx}
\par\smallskip
\normalsize

\noindent\footnotesize\textit{Note.} Each row reports the configuration selected by a four-trial Bayesian sweep that maximized validation macro F1. ``Epochs'' is the training budget selected by the sweep; each complete configuration was subsequently evaluated with seeds 42, 43, and 44. The downstream calibration and routing analyses retained the independently fixed operational DistilBERT checkpoint (learning rate $8.996\times10^{-5}$, batch 32, three epochs, weight decay $10^{-2}$), which achieved 96.69\% held-out accuracy and was not replaced after the matched comparison.\normalsize
\par\medskip



\Needspace{39\baselineskip}
\noindent\textsc{Appendix Table A1d. Phase~1 computational measurements}\par
\label{tab:appendix-a1d-efficiency}
\smallskip
\small
\setlength{\tabcolsep}{4pt}
\renewcommand{\arraystretch}{1.14}
\noindent\textbf{(a) One-off classifier development}\par\smallskip
\noindent\begin{tabularx}{\textwidth}{@{}l *{5}{>{\centering\arraybackslash}X}@{}}
\toprule
Model & Feature extraction (GPU-h) & Search (GPU-h) & Three-seed fitting (GPU-h) & Build total (GPU-h) & Peak memory: extraction / fitting (GB) \\
\midrule
DistilBERT & --- & 0.32 & 0.27 & 0.59 & --- / 4.35 \\
Mistral~7B & 0.67 & 0.02 & 0.003 & 0.70 & 29.11 / 1.11 \\
Llama~2~7B & 0.64 & 0.03 & 0.01 & 0.68 & 27.01 / 1.11 \\
\bottomrule
\end{tabularx}
\par\medskip
\noindent\textbf{(b) Repeated inference measurements}\par\smallskip
\noindent\begin{tabularx}{\textwidth}{@{}l *{5}{>{\centering\arraybackslash}X}@{}}
\toprule
Model & Batched time (min) & Single-post p95 (ms) & Single-post peak memory (GB) & GPU-h per 1,000 posts & Batched GPU energy (Wh) \\
\midrule
DistilBERT & 0.040 & 5.11 $\pm$ 0.12 & 0.16 & 0.000056 & 0.16 \\
Mistral~7B & 3.87 & 38.37 $\pm$ 0.18 & 14.28 & 0.005376 & 25.5 \\
Llama~2~7B & 3.61 & 36.23 $\pm$ 0.40 & 13.27 & 0.005008 & 23.8 \\
\bottomrule
\end{tabularx}
\par\medskip
\noindent\textbf{(c) Complete measured run, including evaluation and repeated benchmarks}\par\smallskip
\noindent\begin{tabularx}{\textwidth}{@{}l *{3}{>{\centering\arraybackslash}X}@{}}
\toprule
Model & Total GPU time (h) & GPU-board energy (Wh) & Batched GPU utilization (\%) \\
\midrule
DistilBERT & 0.61 & 124.6 & 57.0 \\
Mistral~7B & 0.91 & 335.7 & 99.7 \\
Llama~2~7B & 0.87 & 310.8 & 99.7 \\
\bottomrule
\end{tabularx}
\par\smallskip
\noindent\textit{Protocol.} All three configurations were measured on NVIDIA A100-SXM4-40GB accelerators in the same software environment. Training used bf16 and the inference benchmark used fp16. The input-length cap was 256 tokens. Batched inference used batch size 32 on all 12,000 held-out posts; single-post latency used 512 posts at batch size one. Each inference benchmark was repeated three times after three warm-up batches. Table~\ref{tab:phase1-performance}(b) reports mean throughput, mean p50 latency, and their across-repetition SDs; panel (b) above reports mean p95 latency and its SD. Displayed zero SDs result from rounding, not an assumption of zero variation.
\par\smallskip
\noindent\textit{Timing and resource accounting.} CUDA was synchronized before and after each timed region, including host-to-device copies; this is a model-inference benchmark rather than full application latency. Non-padding input tokens were counted. NVML sampled utilization and board power at 1~Hz, and energy was computed by trapezoidal integration of board-power samples. Energy covers the GPU board only; the short DistilBERT timing windows limit the temporal resolution of its energy and utilization estimates. Batched time and energy in panel (b) describe one repetition, not the sum of three. The GPU-hours per 1,000 posts are normalized from batched throughput. Values are rounded independently.
\par\smallskip
\noindent\textit{Build and evaluation scope.} Probe feature extraction processed 108,000 posts before fitting cached-feature linear heads. The three-seed fitting and fitting-memory entries therefore exclude the backbone, whose cost and peak memory appear separately in panel (a). The search used four trials per model. Cached-feature head scoring is not classifier inference throughput; neither those timings nor the DistilBERT training-pipeline scoring rate are used in Table~\ref{tab:phase1-performance}(b). Panel (c) includes search, fitting, scoring, and all timing repetitions, and is not a per-deployment workload estimate. Monetary conversion and Phase~2 generation costs are not included.
\normalsize
\par\medskip

\Needspace{24\baselineskip}
\section{Chain-of-Thought Prompting Protocol}

This table summarizes the final Chain-of-Thought prompting protocol used with Llama~2-7B-Chat to re-evaluate low-confidence Phase~1 predictions. The protocol requires independent text-grounded assessment before Phase~1-label comparison, uses trajectory reasoning only when a shift is supported by the text, and ends with one extractable final label. The earlier numerical emotion-breakdown requirement is not used.

\par\medskip
\noindent\textsc{Appendix Table A2. Chain-of-Thought prompting protocol for emotional re-evaluation}\par
\label{tab:appendix-a2-cot-prompt}
\smallskip
\small
\setlength{\tabcolsep}{3.5pt}
\renewcommand{\arraystretch}{1.2}
\noindent\begin{tabularx}{\textwidth}{@{}>{\raggedright\arraybackslash}p{0.18\textwidth}YYY@{}}
\toprule
Prompt Section & Prompt Instruction & Purpose in the Framework & Expected Output \\
\midrule
Role & You are an expert annotator for research-oriented, non-clinical emotion classification. Do not make a clinical diagnosis, infer a medical condition, or provide treatment advice. & Defines a text-classification role without clinical diagnostic authority. & A non-clinical, research-oriented analysis. \\
Independent assessment & Before seeing the Phase~1 label, identify the dominant emotional meaning of the full text and cite the strongest textual evidence. & Intended to limit anchoring on the initial prediction; this effect was not measured separately. & An initial Depression, Neutral, or Happy assessment. \\
Classification and trajectory rule & Depression requires dominant unresolved distress; Neutral requires mainly factual, routine, balanced, or emotionally mild content; Happy requires dominant happiness, relief, gratitude, accomplishment, fulfillment, or clear positive resolution. Use trajectory reasoning only when a temporal shift is explicitly supported. & Specifies full-text decision criteria and discourages unsupported temporal interpretations. & A text-grounded dominant-emotion decision. \\
Phase~1 comparison & Reveal the Phase~1 label only after independent assessment. Confirm it only when it is supported by the full-text evidence; otherwise explain the correction. & Separates independent assessment from confirm-or-correct reasoning. & A supported confirmation or correction. \\
Output constraint & Use only Depression, Neutral, or Happy. Do not use synonyms or a percentage breakdown. End with exactly one line in the format \texttt{Final label: [label]}. & Produces a stable, machine-extractable final Phase~2 label. & A response ending with one permitted final label. \\
\bottomrule
\end{tabularx}
\par\medskip


\noindent\footnotesize\textit{Note.} This table summarizes the protocol and uses a non-clinical rendering of the role description; it is not a verbatim prompt transcript. The executable prompt strings are retained in the released notebooks. The model's assigned role does not confer diagnostic or treatment authority.\normalsize


\Needspace{34\baselineskip}
\section{SELF-DISCOVER Prompting Protocol}

This table presents the input-specific SELF-DISCOVER adaptation used with Llama~3-8B-Instruct. Discovery is repeated for each routed post rather than once for the whole task as in the original method \cite{ref18}. The protocol requests an emotional assessment, comparison with Phase~1, textual justification, and one final label within the three-class label space.

\par\medskip
\noindent\textsc{Appendix Table A3. SELF-DISCOVER prompting protocol for nuanced emotion classification}\par
\label{tab:appendix-a3-self-discover-prompt}
\smallskip
\footnotesize
\setlength{\tabcolsep}{3pt}
\renewcommand{\arraystretch}{1.18}
\noindent\begin{tabularx}{\textwidth}{@{}>{\raggedright\arraybackslash}p{0.18\textwidth}YYY@{}}
\toprule
Prompt Section & Prompt Instruction & Purpose in the Framework & Expected Output \\
\midrule
Role & You are an expert annotator for mental-health-related emotion classification. Your task is to analyze the emotional content of text data using structured reasoning. This task is intended for research-oriented text classification, not clinical diagnosis. & Assigns the model a structured annotation role for deeper emotional interpretation without implying clinical diagnostic authority. & The model performs structured emotional classification within a non-clinical framework. \\
Task Context & I will provide you with text data and an AI-generated emotional classification label. Your task is to determine the dominant emotion that best represents the overall sentiment of the text. The text may contain multiple emotions, but your goal is to determine the most representative emotion that captures the overall tone. Do not make a clinical diagnosis, infer a medical condition, or provide treatment advice. & Defines the input information, classification objective, and need to handle mixed emotional cues while limiting the task to non-clinical text classification. & The model identifies the dominant emotion while considering the overall emotional trajectory of the text. \\
Question 1: Independent Emotion Analysis & Objectively analyze the given text.  Identify the dominant emotion that best represents the text's overall sentiment. & Encourages independent assessment of the input text before comparing it with the AI-generated label. & An initial dominant emotion based on the text itself. \\
Question 2: Comparison with AI Label & Compare your analysis from Question 1 with the AI's classified label. & Evaluates whether the independent analysis aligns with the Phase 1 prediction. & A comparison between the independent analysis and the AI-generated label. \\
Question 3: Evidence-Based Re-Evaluation & Evaluate the AI's classification using textual evidence. If it aligns with your independent assessment, confirm it. If it does not, determine the correct dominant emotion and justify the decision using only evidence from the text. & Enables reasoning-based correction when the Phase 1 prediction is inconsistent with the textual evidence. & A confirmed or corrected label with a text-grounded justification. \\
Question 4: Final Label Selection & Select the option that most accurately represents your analysis. Do not create or explain additional choices beyond the provided options. End the response with exactly one final label using the format \texttt{Final label: [label]}. Options: Depression, Neutral, Happy. & Restricts the final output to the predefined three-class label space and supports deterministic label extraction. & A structured response ending with one final label selected from Depression, Neutral, and Happy. \\
Classification Guidelines & If the text expresses ongoing sadness, hopelessness, emotional distress, or a strongly negative final emotional trajectory, classify it as Depression. If the text is mainly factual, balanced, informational, or low-affect, classify it as Neutral. If the text expresses happiness, accomplishment, relief, gratitude, or fulfillment as the dominant final sentiment, classify it as Happy. When multiple emotions are present, choose the class that best reflects the overall message and final takeaway rather than isolated phrases; do not default to Neutral merely because cues are mixed. & Provides label-specific decision rules and supports consistent handling of mixed or shifting emotional trajectories. & Consistent application of the three-class emotion-label criteria. \\
Output Constraint & Do not create labels outside the provided options. When explaining the decision, base the justification on textual evidence rather than clinical assumptions. The response should end with exactly one final label in the format Final label: Depression, Final label: Neutral, or Final label: Happy. & Ensures that the output remains aligned with the study's three-class classification setting and can be parsed into a final Phase 2 label. & A structured output that includes a rationale and a clearly extractable final label. \\
\bottomrule
\end{tabularx}
\par\medskip


\noindent\footnotesize\textit{Note.} This table summarizes the protocol with non-clinical role wording; verbatim executable prompts are retained in the released notebooks.\normalsize

\par\medskip
\Needspace{30\baselineskip}
\begin{center}
\includegraphics[width=0.82\textwidth]{figures/appendix_c_self_discover_workflow_vector_final.pdf}
\end{center}
\vspace{-0.5\baselineskip}
\noindent\footnotesize\textbf{Appendix Figure C1.} Per-input SELECT--ADAPT--IMPLEMENT plan construction in the study's SELF-DISCOVER adaptation.\normalsize\par
\smallskip
\noindent\footnotesize\textit{Note.} The diagram depicts the study implementation, not a new reasoning method. Unlike the task-level plan reuse in \cite{ref18}, a plan is constructed for each routed post. Modules are prompt instructions, not trained neural components. The plan is retained for inspection; quantitative evaluation uses the terminal label.\normalsize
\par\medskip
\Needspace{20\baselineskip}
\noindent\textsc{Appendix Table A3b. Canonical 39-module pool used by SELF-DISCOVER}\par
\label{tab:appendix-a3b-module-pool}
\smallskip
\footnotesize
\setlength{\tabcolsep}{3.5pt}
\renewcommand{\arraystretch}{1.14}
\noindent\begin{tabularx}{\textwidth}{@{}>{\raggedright\arraybackslash}p{0.13\textwidth}>{\raggedright\arraybackslash}p{0.25\textwidth}YY@{}}
\toprule
Module ID(s) & Canonical family or named style & Function in the supplied pool & Relevance to routed emotion re-evaluation \\
\midrule
1--9 & General problem framing & Experiment design, candidate ideas, progress measurement, simplification, assumptions, risks, alternative perspectives, long-term implications, and decomposition. & Offers generic operations from which the model may select a smaller task-relevant subset. \\
10 & Critical Thinking & Examine evidence from multiple perspectives, question assumptions, and identify bias or logical weaknesses. & Supports comparison of isolated emotional cues with evidence from the complete post. \\
11 & Creative Thinking & Generate unconventional alternatives beyond a default solution. & Provides alternative interpretations when the initial label is poorly supported. \\
12 & Collaborative Thinking & Seek diverse perspectives and external expertise. & Included in the canonical pool, although the single-model implementation does not contact external reviewers. \\
13 & Systems Thinking & Examine interdependencies, underlying causes, feedback, and the larger system. & Can connect local statements with the broader narrative context of the post. \\
14 & Risk Analysis & Evaluate uncertainty, trade-offs, and consequences of competing choices. & Can prompt caution when evidence for the candidate labels is incomplete or conflicting. \\
15 & Reflective Thinking & Reconsider assumptions, biases, and prior interpretations. & Supports re-examination of the Phase~1 prediction after an independent reading. \\
16--24 & Diagnosis and evidence planning & Identify the core issue, causes, prior attempts, obstacles, data, affected stakeholders, resources, progress measures, and metrics. & Supplies problem-diagnosis and evidence-checking operations that can be adapted to emotional text. \\
25--32 & Problem-type and constraint checks & Determine whether a problem is technical, physical, behavioral, decision-based, analytical, design-oriented, systemic, or urgent. & Helps the model decide which reasoning style is relevant; these prompts are not clinical assessments. \\
33--37 & Solution-family and counterfactual revision & Recall typical solutions, propose alternatives, challenge the current best solution, modify it, or replace it. & Encourages alternatives to simply accepting the Phase~1 label. \\
38--39 & Step-by-step planning and execution & Construct and implement an explicit sequential plan. & Supports the structured \texttt{IMPLEMENT} artifact used before the final answer. \\
\bottomrule
\end{tabularx}
\par\smallskip
\noindent\footnotesize\textit{Note.} These are prompt-level reasoning instructions, not learned model components. The table groups the verbatim 39-entry pool for readability; the complete wording and numbering are embedded in the released final Colab notebooks. For each routed post, \texttt{SELECT} is model-generated and may paraphrase or group canonical entries. The study therefore audits the trace but scores only the parsed final class label.\normalsize

\par\medskip
\Needspace{18\baselineskip}
\noindent\textsc{Appendix Table A3c. Exploratory cleaned-input prompt-policy diagnostics}\par
\label{tab:appendix-a3c-prompt-ablation}
\smallskip
\footnotesize
\setlength{\tabcolsep}{3pt}
\renewcommand{\arraystretch}{1.16}
\noindent\begin{tabularx}{\textwidth}{@{}>{\raggedright\arraybackslash}p{0.20\textwidth}>{\raggedright\arraybackslash}p{0.28\textwidth}Y>{\raggedright\arraybackslash}p{0.15\textwidth}@{}}
\toprule
Prompt policy & Defining instruction & Routed-only development result & Decision \\
\midrule
Llama~2 final CoT prompt & Analyze the text independently, compare with the Phase~1 label only afterward, and end with one exact class label. & 42.69\%; 38 corrected and 49 introduced (net $-11$); no parse failure. & Retained for final rerun \\
Llama~3 shared cross-model prompt & Apply the same trajectory-aware instruction template considered for both LLM families. & 43.90\%; 32 introduced errors; seven terminal labels were not parsed reliably. & Rejected \\
Llama~3 error-aware refinement & Add explicit evidence checks and a stricter terminal-label contract to the shared prompt. & 50.29\%; parsing was repaired, but routed-only accuracy remained below SELF-DISCOVER. & Rejected \\
Llama~3 SELF-DISCOVER & Dynamically SELECT, ADAPT, and IMPLEMENT reasoning modules before producing one final class label. & 52.63\%; 13 corrected and seven introduced (net $+6$); no parse failure. & Retained for final rerun \\
\bottomrule
\end{tabularx}
\par\smallskip
\noindent\footnotesize\textit{Note.} These are exploratory development diagnostics from the earlier cleaned-input condition, not the final original-text end-to-end result. Internal prompt version identifiers are omitted because they are experiment-management labels rather than published methods. All rows used the same Phase~1 model, calibrated temperature, routing threshold, 171 routed IDs, LLM checkpoint family, generation settings, and evaluation definitions. The final original-text results are reported in Tables~\ref{tab:reddit-e2e-results} and~\ref{tab:paired-statistics}.\normalsize
\par\medskip
\normalsize


\Needspace{12\baselineskip}
\section{Observed Routing and Reasoning Audit Trail}

This appendix retains one observed Mixed Emotion routed correction for each Phase~2 model and one observed Reddit original-text correction from the final Llama~3 run. The tables present stored output fields rather than reconstructed narratives, so a reader can trace the original Phase~1 label, text-grounded reasoning, terminal label, and evaluation reference. These are research artifacts for proxy emotion classification, not clinical diagnostic reports.

For every routed Llama~3 example, the result CSV retains \texttt{LLaMA3\_SELECT}, \texttt{LLaMA3\_ADAPT}, \texttt{LLaMA3\_IMPLEMENT}, \texttt{LLaMA3\_Answer}, and the parsed \texttt{LLaMA3\_final\_label}; analogous Llama~2 records retain its independent analysis, comparison, and final-response fields. Thus, aggregate results can be traced back to case-level reasoning artifacts without treating generated text as clinical evidence.

\par\medskip
\Needspace{14\baselineskip}
\noindent\textsc{Appendix Table A4a. Llama~2 CoT output-column example for an observed correction}\par
\label{tab:appendix-a4-llama2-output-columns}
\smallskip
\footnotesize
\setlength{\tabcolsep}{3.5pt}
\renewcommand{\arraystretch}{1.16}
\noindent\begin{tabularx}{\textwidth}{@{}>{\raggedright\arraybackslash}p{0.18\textwidth}>{\raggedright\arraybackslash}p{0.43\textwidth}Y@{}}
\toprule
Stored column & Observed output excerpt for MEV2\_DEP\_062 & Function and interpretation \\
\midrule
Input and reference & ``At first, a workday felt manageable because there was encouraging feedback from someone I respect. Still, I kept track of dates, forms, and reminders, and as the day went on, there was a sharp drop in my mood once I had time to think. I tried to stay honest about the mixed nature of the experience, especially before going to sleep. What stayed longest was not the good moment, but the sadness that returned and remained unresolved.'' Reference: Depression; Phase~1: Happy. & Establishes the exact evidence, reference label, and initial error being audited. The Phase~2 model must correct the label from the text rather than from hidden metadata. \\
\texttt{LLaMA2\_1} & ``The dominant emotion is sadness.'' The output cites the encouraging feedback, subsequent sharp mood drop, and sadness that remained unresolved, concluding that the negative trajectory dominates the initial positive cue. & Records an assessment before Phase~1 comparison and its stated emphasis on the final trajectory; it does not establish the cause of the prediction. \\
\texttt{LLaMA2\_2} & ``I have to disagree with the Phase 1 classifier's prediction of Happy.'' The output again cites the mood drop and unresolved sadness as evidence that the initial label is inconsistent with the complete text. & Audits agreement with Phase~1 by naming the disputed label and the text-grounded reason for rejecting it. \\
\texttt{LLaMA2\_3} & \texttt{Final label: Depression} & Supplies the canonical terminal response. The parser accepts only Depression, Neutral, or Happy and uses this label to replace the routed Phase~1 prediction. \\
\texttt{LLaMA2\_final\_label} & Depression & Stores the machine-extractable result used in end-to-end metrics and the confusion matrix; free-form prose is retained only for audit. \\
\bottomrule
\end{tabularx}
\par\medskip
\normalsize

\Needspace{26\baselineskip}
\noindent\textsc{Appendix Table A4b. Llama~3 SELF-DISCOVER output-column example for an observed correction}\par
\label{tab:appendix-a4-llama3-output-columns}
\smallskip
\footnotesize
\setlength{\tabcolsep}{3.5pt}
\renewcommand{\arraystretch}{1.16}
\noindent\begin{tabularx}{\textwidth}{@{}>{\raggedright\arraybackslash}p{0.18\textwidth}>{\raggedright\arraybackslash}p{0.43\textwidth}Y@{}}
\toprule
Stored column & Observed output excerpt for MEV2\_DEP\_022 & Function and interpretation \\
\midrule
Input and reference & ``At first, a workday felt manageable because there was encouraging feedback from someone I respect. Still, I kept track of dates, forms, and reminders, and as the day went on, there was a sharp drop in my mood once I had time to think. I realized that the ending of the day mattered more than the beginning, especially during a slow part of the evening. What stayed longest was not the good moment, but the sadness that returned and remained unresolved.'' Reference: Depression; Phase~1: Happy. & Establishes the original mixed-trajectory input and the Phase~1 error. It lets the reader compare the later reasoning trace directly with the evidence available at inference time. \\
\texttt{LLaMA3\_SELECT} & The stored trace selects Critical Thinking, Textual Analysis, Risk Analysis, and Reflective Thinking. Its preliminary analysis notes that the post begins positively but ends with unresolved sadness and emotional distress. & Records the actual model-generated selection trace. The generated names may paraphrase entries in the canonical pool and are retained for inspection rather than treated as scored predictions. \\
\texttt{LLaMA3\_ADAPT} & The selected operations are rewritten as Emotion-Focused Analysis, Textual Evidence-Based Judgment, Emotion Classification Risk Assessment, and Reflective Emotional Understanding. & Tailors the selected operations to the three-class task while requiring the decision to consider the full trajectory and alternative interpretations. \\
\texttt{LLaMA3\_IMPLEMENT} & The JSON plan defines the three class guidelines, then sequences emotional-cue identification, textual-evidence review, classification-risk assessment, reflective review, and terminal-label production. & Records the generated plan before execution; the plan can be inspected, but only the final parsed label contributes to accuracy. \\
\texttt{LLaMA3\_Answer} & ``The text starts with a positive note ... but this is short-lived. The speaker's mood drops sharply ... and the text concludes with unresolved sadness.'' It rejects Happy and ends with \texttt{Final label: Depression}. & Provides the human-readable, text-grounded assessment and directly explains why the initial Happy prediction is corrected. \\
\texttt{LLaMA3\_final\_label} & Depression & Stores the parsed terminal label used only for routed examples; accepted examples retain their Phase~1 label. \\
\bottomrule
\end{tabularx}
\par\medskip
\normalsize

Note. Tables~A4a and A4b are publication-readable renderings of actual stored routed-sample outputs. Output excerpts are shortened only for layout; the full text fields remain in the Llama~2 and Llama~3 result CSVs. The observed Mixed Emotion results are not clinical diagnostic reports.

\noindent\textsc{Appendix Table A4c. Mixed Emotion end-to-end accuracy by controlled scenario}\par
\label{tab:appendix-a4c-scenario-results}
\smallskip
\footnotesize
\setlength{\tabcolsep}{3.5pt}
\renewcommand{\arraystretch}{1.14}
\noindent\begin{tabularx}{\textwidth}{@{}>{\raggedright\arraybackslash}X c c c c c@{}}
\toprule
Controlled scenario & $n$ & Routed & Phase~1 & Llama~2 e2e & Llama~3 e2e \\
\midrule
Blended-emotion co-occurrence & 40 & 8 & 60.0\% & 70.0\% & 72.5\% \\
Positive-to-distress shift & 40 & 9 & 65.0\% & 77.5\% & 77.5\% \\
Distress-to-recovery shift & 40 & 12 & 90.0\% & 95.0\% & 97.5\% \\
Neutral framing with subtle affect & 40 & 10 & 87.5\% & 92.5\% & 97.5\% \\
Conflicting cues, dominant trajectory & 40 & 0 & 100.0\% & 100.0\% & 100.0\% \\
Clear technical/informational Neutral & 75 & 0 & 100.0\% & 100.0\% & 100.0\% \\
Mild factual/procedural Neutral ambiguity & 25 & 5 & 32.0\% & 28.0\% & 36.0\% \\
\bottomrule
\end{tabularx}
\par\smallskip
\noindent\footnotesize\textit{Note.} Accepted examples retain their Phase~1 labels, so rows with zero routed examples have identical end-to-end accuracy. These scenario results are descriptive because the controlled groups are small and unequal in size. They show where the aggregate gain originates without treating the synthetic scenario distribution as representative of real-world prevalence.\normalsize

\par\medskip
\noindent\textsc{Appendix Table A4d. Observed Reddit original-text Llama~3 correction record}\par
\label{tab:appendix-a4d-reddit-original-correction}
\smallskip
\footnotesize
\setlength{\tabcolsep}{3.5pt}
\renewcommand{\arraystretch}{1.15}
\noindent\begin{tabularx}{\textwidth}{@{}>{\raggedright\arraybackslash}p{0.19\textwidth}Y@{}}
\toprule
Stored field & Observed record for RED\_077578 \\
\midrule
Original Phase~2 input & ``GF is `numb' and is tired of trying to stay positive. What can I do?'' The stored proxy reference is Depression, while Phase~1 predicted Happy. \\
\texttt{LLaMA3\_SELECT} & Selected evidence review and reflective operations; the trace identified emotional numbness, exhaustion, and the absence of a positive resolution before comparing with the Happy prediction. \\
\texttt{LLaMA3\_ADAPT} & Reframed the selected operations as emotion extraction, analysis, comparison, reconciliation, and three-class classification. \\
\texttt{LLaMA3\_IMPLEMENT} & Generated a JSON plan that identifies emotional cues, evaluates the overall trajectory, compares the independent assessment with Phase~1, and produces one terminal label. \\
\texttt{LLaMA3\_Answer} & ``The text expresses feelings of being `numb' and being tired of trying to stay positive, which indicates a sense of emotional exhaustion and hopelessness.'' It rejects Happy and ends with \texttt{Final label: Depression}. \\
Evaluation record & Phase~1: Happy (incorrect); Phase~2: Depression (correct); counted as one corrected routed error in the final original-text Reddit run. \\
\bottomrule
\end{tabularx}
\par\smallskip
\noindent\footnotesize\textit{Note.} This is an observed stored output, shortened only for publication layout. URLs and direct username patterns were masked before inference; no such pattern occurred in this excerpt. The rationale is an audit artifact and is not treated as a clinical explanation.\normalsize

\par\medskip
\Needspace{22\baselineskip}
\noindent\textsc{Appendix Table A4e. Audit classification of representative accepted high-confidence disagreements}\par
\label{tab:appendix-a4e-accepted-errors}
\smallskip
\footnotesize
\setlength{\tabcolsep}{3.2pt}
\renewcommand{\arraystretch}{1.16}
\noindent\begin{tabularx}{\textwidth}{@{}p{0.10\textwidth}p{0.15\textwidth}c p{0.24\textwidth}Y@{}}
\toprule
ID & Reference $\rightarrow$ prediction & Conf. & Observed failure mode & Evidence in the original post \\
\midrule
RED\_041439 & Depression $\rightarrow$ Happy & 0.983 & Acute-distress cue missed & The title explicitly states renewed suicidal ideation; the preceding ``okay'' period is a possible conflicting cue, not a verified cause of the prediction. \\
RED\_000879 & Depression $\rightarrow$ Happy & 0.980 & Negative-state cue missed & Blackmail, alcohol dependence, imprisonment, family separation, and loss of custody provide no textual support for Happy. \\
RED\_026909 & Happy $\rightarrow$ Neutral & 0.995 & Technical-context dominance & A programming narrative ends with explicit unexpected achievement and pride, yet the prediction is Neutral; the role of the technical vocabulary remains a hypothesis. \\
RED\_103680 & Happy $\rightarrow$ Neutral & 0.989 & Factual-detail dominance & Explicit excitement about a successful purchase and new project is obscured by product and pricing details. \\
RED\_010855 & Neutral $\rightarrow$ Depression & 0.982 & Topic-term shortcut & A factual question about computerized diagnosis is classified as Depression after repeated mental-health terminology. \\
RED\_000438 & Neutral $\rightarrow$ Happy & 0.986 & Question-form/context error & A factual question about a sound from a cup of tea contains no positive affect but receives a highly confident Happy prediction. \\
\bottomrule
\end{tabularx}
\par\smallskip
\noindent\footnotesize\textit{Note.} All six rows are accepted errors because prediction $\ne$ stored reference proxy label and confidence $\geq0.70$. Cases were selected only after exact original-post linkage and textual review; they do not estimate failure-mode prevalence.\normalsize

\par\medskip
\Needspace{20\baselineskip}
\noindent\textsc{Appendix Table A4f. Verifiable case records for the accepted high-confidence audit}\par
\label{tab:appendix-a4f-accepted-error-records}
\smallskip
\noindent\footnotesize Each excerpt was recovered from the retained \texttt{title} and \texttt{selftext} fields by normalized exact matching to the stored model input. Capitalization, punctuation, and discourse order are preserved; URLs and unnecessary identifying details are omitted. The classifier itself received the corresponding cleaned field.\normalsize

\par\smallskip
\noindent\begin{tabularx}{\textwidth}{@{}>{\raggedright\arraybackslash}p{0.18\textwidth}Y@{}}
\toprule
\multicolumn{2}{@{}l}{\textbf{RED\_041439 --- explicit acute distress missed}} \\
\midrule
Original-post excerpt & ``Aaaannd I'm suicidal again... The past two weeks were going pretty okay too....'' \\
Recorded decision & Reference proxy: Depression; Phase~1 prediction: Happy; calibrated confidence: 0.983; accepted at $\tau=0.70$. \\
Text-grounded audit & The present-state title directly expresses suicidal ideation. The statement that the preceding two weeks were ``okay'' does not reverse the current distress, so Happy is inconsistent with the post's final state. \\
Evaluation treatment & Counted as one accepted Depression $\rightarrow$ Happy error; no Phase~2 output exists because the row was not routed. \\
\bottomrule
\end{tabularx}

\par\medskip
\noindent\begin{tabularx}{\textwidth}{@{}>{\raggedright\arraybackslash}p{0.18\textwidth}Y@{}}
\toprule
\multicolumn{2}{@{}l}{\textbf{RED\_000879 --- severe negative state predicted as Happy}} \\
\midrule
Original-post excerpt & ``My ex wife blackmailed me into signing child support paperwork and made me leave my home. Today I am an alcoholic headed to prison and [...] she has custody of my little boy.'' \\
Recorded decision & Reference proxy: Depression; Phase~1 prediction: Happy; calibrated confidence: 0.980; accepted at $\tau=0.70$. \\
Text-grounded audit & The post describes coercion, alcohol dependence, imprisonment, separation from home, and loss of custody. No positive resolution or Happy cue appears in the source text. \\
Evaluation treatment & Counted as one accepted Depression $\rightarrow$ Happy error under the unchanged reference-label evaluation. \\
\bottomrule
\end{tabularx}

\par\medskip
\noindent\begin{tabularx}{\textwidth}{@{}>{\raggedright\arraybackslash}p{0.18\textwidth}Y@{}}
\toprule
\multicolumn{2}{@{}l}{\textbf{RED\_026909 --- explicit pride masked by technical context}} \\
\midrule
Original-post excerpt & ``I'm Teaching Myself Java. My husband and I are trying to learn how to program [...] I never expected to grasp it as quickly as I am, and I'm kinda proud of myself.'' \\
Recorded decision & Reference proxy: Happy; Phase~1 prediction: Neutral; calibrated confidence: 0.995; accepted at $\tau=0.70$. \\
Text-grounded audit & Programming details dominate most of the post, but the explicitly stated takeaway is unexpected accomplishment and pride. Neutral misses the positive evaluative conclusion. \\
Evaluation treatment & Counted as one accepted Happy $\rightarrow$ Neutral error. \\
\bottomrule
\end{tabularx}

\clearpage
\noindent\textsc{Appendix Table A4f (continued). Verifiable case records for the accepted high-confidence audit}\par
\smallskip
\noindent\begin{tabularx}{\textwidth}{@{}>{\raggedright\arraybackslash}p{0.18\textwidth}Y@{}}
\toprule
\multicolumn{2}{@{}l}{\textbf{RED\_103680 --- excitement obscured by product details}} \\
\midrule
Original-post excerpt & ``I just saved \$170 off the cost of an EEPROM programmer by shopping on eBay---my new project/hobby couldn't start off with a better investment! I'm so excited! [...] Compare \$28 [...] to \$200 plus shipping.'' \\
Recorded decision & Reference proxy: Happy; Phase~1 prediction: Neutral; calibrated confidence: 0.989; accepted at $\tau=0.70$. \\
Text-grounded audit & The pricing and hardware terminology are factual, but the opening and stated evaluation are explicitly enthusiastic. The error is consistent with topic sensitivity, but the output alone does not establish the classifier's cue attribution. \\
Evaluation treatment & Counted as one accepted Happy $\rightarrow$ Neutral error. \\
\bottomrule
\end{tabularx}

\par\medskip
\noindent\begin{tabularx}{\textwidth}{@{}>{\raggedright\arraybackslash}p{0.18\textwidth}Y@{}}
\toprule
\multicolumn{2}{@{}l}{\textbf{RED\_010855 --- mental-health terminology in a factual question}} \\
\midrule
Original-post excerpt & ``Sorry if this is a common question but, why can't computers diagnose mental health disorders? [...] why can't an online test diagnose it? Is it just the authority they have to officially diagnose things and prescribe medicine and treatment?'' \\
Recorded decision & Reference proxy: Neutral; Phase~1 prediction: Depression; calibrated confidence: 0.982; accepted at $\tau=0.70$. \\
Text-grounded audit & The post discusses depression-related concepts but does not express the author's emotional state. Its function is an informational question, consistent with a possible topic-vocabulary shortcut, rather than demonstrating that mechanism. \\
Evaluation treatment & Counted as one accepted Neutral $\rightarrow$ Depression error. \\
\bottomrule
\end{tabularx}

\noindent\begin{tabularx}{\textwidth}{@{}>{\raggedright\arraybackslash}p{0.18\textwidth}Y@{}}
\toprule
\multicolumn{2}{@{}l}{\textbf{RED\_000438 --- affect inferred from a factual science question}} \\
\midrule
Original-post excerpt & ``Why do I hear a high-pitched tone coming from my tea? I boiled some water, and poured it onto a tea bag, inside of a ceramic mug. [...] What's going on? And is it safe to drink?'' \\
Recorded decision & Reference proxy: Neutral; Phase~1 prediction: Happy; calibrated confidence: 0.986; accepted at $\tau=0.70$. \\
Text-grounded audit & The post reports an observation and asks for an explanation and safety information. It contains no explicit positive emotional evaluation, making Happy unsupported by the original text. \\
Evaluation treatment & Counted as one accepted Neutral $\rightarrow$ Happy error. \\
\bottomrule
\end{tabularx}

\par\smallskip
\noindent\footnotesize\textit{Audit scope.} These six cases illustrate error patterns, not verified causal mechanisms or their prevalence. Quotations retain source wording except for URL removal and bracketed shortening. The complete prediction-level accepted-error file is included with the reproducibility package.\normalsize

\par\medskip
\noindent\textsc{Appendix Table A4g. Routing concentration and conditional correction opportunity}\par
\label{tab:appendix-a4g-routing-decomposition}
\smallskip
\footnotesize
For an evaluation set of size $N$ with $E$ Phase~1 errors, let $n_R$ be the number routed and $E_R$ the number of Phase~1 errors among them. The routed-error enrichment and conditional oracle accuracy are
\[
\mathrm{Enrichment}=\frac{E_R/n_R}{E/N},
\qquad
\mathrm{Acc}_{\mathrm{oracle}\mid R}=\mathrm{Acc}_{\mathrm{P1}}+\frac{E_R}{N}.
\]
Accuracies in the equation are proportions; the table reports percentages. For reasoner $m$, the realized share of the fixed routing opportunity is
\[
\rho_m=\frac{\mathrm{Corrected}_m-\mathrm{Introduced}_m}{E_R}.
\]
This share is negative when introduced errors outnumber corrections.
\par\smallskip
\footnotesize
\setlength{\tabcolsep}{1.5pt}
\renewcommand{\arraystretch}{1.15}
\noindent\begin{tabularx}{0.98\textwidth}{@{}l c c c c c c c@{}}
\toprule
Dataset & Full err. & Routed err. & Enrichment & Routed P1 & Llama~2 routed & Llama~3 routed & Oracle e2e \\
\midrule
Reddit & 3.31\% & 50.88\% & 15.38$\times$ & 49.12\% & 47.37\% & 66.67\% & 97.42\% \\
Mixed Emotion & 18.67\% & 47.73\% & 2.56$\times$ & 52.27\% & 79.55\% & 93.18\% & 88.33\% \\
\bottomrule
\end{tabularx}
\par\smallskip
\noindent\footnotesize\textit{Note.} Reddit has $N=12{,}000$, $E=397$, $n_R=171$, and $E_R=87$; Mixed Emotion has $N=300$, $E=56$, $n_R=44$, and $E_R=21$. Under the fixed router, Llama~2 and Llama~3 realize $-3.45\%$ and 34.48\% of the Reddit correction opportunity, and 57.14\% and 85.71\% of the Mixed Emotion opportunity. The oracle assumes perfect correction of routed Phase~1 errors with no introduced errors. It is a descriptive upper bound under the fixed routing policy, not a trained comparator and not a basis for threshold selection.\normalsize


\par\medskip
\Needspace{16\baselineskip}
\section{Synthetic Mixed Emotion Dataset Generation Protocol}

\par\medskip
\noindent\textsc{Appendix Table A5. Synthetic Mixed Emotion Dataset generation protocol}\par
\label{tab:appendix-a5-mixed-emotion-protocol}
\smallskip
\footnotesize
\setlength{\tabcolsep}{4pt}
\renewcommand{\arraystretch}{1.08}
\noindent\begin{tabularx}{\textwidth}{@{}>{\raggedright\arraybackslash}p{0.20\textwidth}Y@{}}
\toprule
Component & Description \\
\midrule
Purpose & Construct 300 controlled examples spanning mixed or shifting affect and Neutral controls for supplementary stress-test evaluation. \\
Use in pipeline & Not used for Phase 1 training, hyperparameter tuning, or routing-threshold selection; used only for supplementary robustness evaluation. \\
Class balance & Depression = 100, Neutral = 100, Happy = 100. \\
Scenario balance & Seven controlled scenario groups: five 40-example mixed or shifting-emotion groups, a 75-example clear technical/informational Neutral group, and a 25-example mild factual/procedural-ambiguity Neutral group. \\
Scenario types & Blended-emotion co-occurrence; positive-to-distress shift; distress-to-recovery shift; Neutral framing with subtle affect; conflicting cues with a dominant trajectory; clear technical or informational Neutral content; and mild factual or procedural Neutral ambiguity. \\
Label rule & Assign Depression, Neutral, or Happy according to the dominant overall emotional trajectory, not isolated phrases. \\
Exclusion criteria & Exclude explicit diagnosis claims, treatment advice, medication references, crisis language, identifying information, and off-scenario content. Mixed-cue requirements do not apply to the clear Neutral controls. \\
Prompt disclosure & Generation instructions are reproduced below; the version, model, date, and inclusion criteria are recorded here. \\
Generation record & GPT-5; prompt version \texttt{mixed-emotion-stress-test-v2.4-neutral-clear}; final records dated Aug. 7, 2026. \\
Limitations & Synthetic origin, generation-model style bias, and lack of expert annotation. \\
\bottomrule
\end{tabularx}
\par\medskip


\noindent\footnotesize\textit{Generation prompt used for synthetic stress-test examples.}\normalsize\par
\smallskip
\noindent
\begin{minipage}[t]{0.49\textwidth}\scriptsize
Generate synthetic Reddit-style posts for a controlled mixed-emotion stress-test dataset. Use one of three target labels: Depression, Neutral, or Happy. For Depression and Happy examples, include mixed or shifting cues but make the final emotional trajectory clear to a trained reviewer. For Neutral examples, retain low affective intensity: include either technical/informational content or mild factual/procedural ambiguity without sustained distress, relief, accomplishment, or other dominant affect. When multiple emotions occur, the target label must follow the overall message and final takeaway rather than an isolated phrase.
\end{minipage}\hfill
\begin{minipage}[t]{0.49\textwidth}\scriptsize
For each example, write one realistic post between 60 and 90 words. Do not include explicit diagnosis claims, medication names, therapy claims, suicide, self-harm, crisis language, usernames, URLs, subreddit names, hashtags, or personally identifying information. Avoid duplicated phrasing and avoid making the label obvious through a single keyword. Return structured records with: example identifier, target label, scenario type, primary context, dominant emotional trajectory, text, and brief label rationale.
\end{minipage}

\par\medskip
\Needspace{19\baselineskip}
\section{Additional Holdout Class-Wise Results}
\label{app:additional-holdout}

Table~A6 supplements the full-set results in Table~\ref{tab:additional-holdout} with class-wise F1 and accuracy on the shared routed subset. Both configurations improved Depression and Happy F1, whereas Neutral F1 decreased slightly. The evaluation protocol is described in Section~\ref{sec:additional-holdout-protocol}.

\par\medskip
\Needspace{13\baselineskip}
\noindent\textsc{Appendix Table A6. Additional holdout: class-wise F1 and routed accuracy}\par
\smallskip
\small
\setlength{\tabcolsep}{5pt}
\renewcommand{\arraystretch}{1.16}
\noindent\begin{tabularx}{\textwidth}{@{}Y c c c c@{}}
\toprule
Configuration & Depression F1 & Neutral F1 & Happy F1 & Routed accuracy \\
\midrule
Phase~1 & 95.87\% & 98.82\% & 95.52\% & 51.09\% \\
+ Llama~2 CoT & 96.23\% & 98.53\% & 95.94\% & 62.04\% \\
+ Llama~3 SELF-DISCOVER & 96.36\% & 98.67\% & 96.07\% & 70.80\% \\
\bottomrule
\end{tabularx}
\par\smallskip
\noindent\footnotesize\textit{Note.} Class F1 uses all 9,000 posts; routed accuracy uses the same 137 posts for each configuration. Both re-evaluators increased overall macro F1 but slightly reduced Neutral F1. These are fixed-run results, not means across training seeds.\normalsize

% REQUIRED BEFORE SUBMISSION:
% 1. Add the final Acknowledgment text, including funding and the IEEE-required
%    disclosure of any AI-generated manuscript text and the affected sections.
% 2. Add one IEEEbiographynophoto or IEEEbiography environment for every author.
% Example:
% \begin{IEEEbiographynophoto}{Full Author Name}
% Short biography required by IEEE Access.
% \end{IEEEbiographynophoto}
\EOD
\end{document}
